%! BibTeX Compiler = biber
%DIF LATEXDIFF DIFFERENCE FILE
%DIF DEL main-first-submission.tex   Wed Feb  4 16:30:12 2026
%DIF ADD main-R1.tex                 Wed Feb  4 16:51:22 2026
\documentclass{article}
\usepackage{caption}
\usepackage{censor}
\usepackage{xcolor, colortbl}
\definecolor{BLUELINK}{HTML}{0645AD}
\definecolor{DARKBLUELINK}{HTML}{0B0080}
\definecolor{LIGHTGREY}{gray}{0.9}
\PassOptionsToPackage{hyphens}{url}
\usepackage[colorlinks=false]{hyperref}
% for linking between references, figures, TOC, etc in the pdf document
\hypersetup{colorlinks,
linkcolor=DARKBLUELINK,
anchorcolor=DARKBLUELINK,
citecolor=DARKBLUELINK,
filecolor=DARKBLUELINK,
menucolor=DARKBLUELINK,
urlcolor=BLUELINK
} % Color citation links in purple
\PassOptionsToPackage{unicode}{hyperref}
\PassOptionsToPackage{naturalnames}{hyperref}

\usepackage{biorxiv}
\usepackage[
style=authoryear,
doi=false,
isbn=false,
url=false,
date=year,
natbib=true,
hyperref,
mincitenames = 1,
maxcitenames = 2,
minbibnames = 1,
maxbibnames = 30,
uniquename=false,
uniquelist=false,
giveninits=true,
backend=biber]{biblatex}
\AtEveryBibitem{\clearfield{note}}
\addbibresource{references_bibtex.bib}

\usepackage{url}
\usepackage{amssymb,amsfonts,amsmath,amsthm,mathtools}
\usepackage{lmodern}
\usepackage{xfrac, nicefrac}
\usepackage{bm}
\usepackage{listings, enumerate, enumitem}
\usepackage[export]{adjustbox}
\usepackage{graphicx}
\usepackage{bbold}
\usepackage{pdfpages}
\pdfinclusioncopyfonts=1
\usepackage{lineno}
\usepackage{tabu}
\usepackage{hhline}
\usepackage{multicol,multirow,array}
\usepackage{etoolbox}
\usepackage{booktabs}
\usepackage{makecell}
\usepackage{marvosym}
\usepackage{orcidlink}

% -- Defining colors:
\definecolor{backcolour}{rgb}{0.95,0.95,0.92}% Definig a custom style:
\lstdefinestyle{mystyle}{
backgroundcolor=\color{backcolour},
basicstyle=\ttfamily\scriptsize\bfseries,
breakatwhitespace=false,
breaklines=true,
captionpos=t,
keepspaces=true,
showspaces=false,
showstringspaces=false,
showtabs=false,
tabsize=2
}% -- Setting up the custom style:
\lstset{style=mystyle}
\captionsetup[table]{hypcap=false}
\captionsetup[figure]{hypcap=false}

\newcommand{\Multiply}{\cdot}
%DIF 83a83-89
\newcommand{\MultiplyMatrix}{\times} %DIF > 
\newcommand{\UniDimArray}[1]{\bm{#1}} %DIF > 
\newcommand{\BiDimArray}[1]{\bm{#1}} %DIF > 
\newcommand{\tr}{^{\intercal}} %DIF > 
\newcommand{\inv}{^{-1}} %DIF > 
\newcommand{\Det}{\text{Det}} %DIF > 
\DeclareMathOperator{\E}{\mathbb{E}} %DIF > 
%DIF -------
\DeclareMathOperator{\Var}{\text{var}}
%DIF 84a91
\DeclareMathOperator{\Cov}{\text{cov}} %DIF > 
%DIF -------
\newcommand{\der}{\mathrm{d}}
\newcommand{\e}{\text{e}}
\newcommand{\Ne}{N_{\text{e}}}
\newcommand{\PhyloSupport}{\pi}
\newcommand{\MeanPhyloSupport}{\overline{\PhyloSupport}}
\newcommand{\Indiv}{k}
\newcommand{\Trait}{P}
\newcommand{\Heritability}{h^2}
\newcommand{\MutationRatePheno}{\mu}
\newcommand{\MutationRateNuc}{u}
\newcommand{\NbrLoci}{L}
\newcommand{\VarEnv}{V_{\mathrm{E}}}
\newcommand{\brownian}{\mathcal{B}}
%DIF 97a105-135
\newcommand{\proba}{\mathbb{P}} %DIF > 
\newcommand{\pfix}{\proba_{\text{fix}}} %DIF > 
\newcommand{\Spi}{i} %DIF > 
\newcommand{\Spj}{j} %DIF > 
\newcommand{\NbrTaxa}{n} %DIF > 
\newcommand{\Time}{T} %DIF > 
\newcommand{\GenerationTime}{\tau} %DIF > 
\newcommand{\NbrGen}{g_{\Spi, \Spj}} %DIF > 
\newcommand{\NucDiv}{d_{\Spi, \Spj}} %DIF > 
\newcommand{\VarGenetic}{V_{\mathrm{A}}} %DIF > 
\newcommand{\HeritabilitySpi}{\Heritability_{\Spi}} %DIF > 
\newcommand{\MeanTrait}{\bar{\Trait}} %DIF > 
\newcommand{\VecTrait}{\UniDimArray{\bar{\Trait}}} %DIF > 
\newcommand{\VarPhy}{\Cov \left( \MeanTrait_{\Spi}, \MeanTrait_{\Spj}\right)} %DIF > 
\newcommand{\VecZero}{\UniDimArray{0}} %DIF > 
\newcommand{\VecOne}{\UniDimArray{1}} %DIF > 
\newcommand{\Distance}{\BiDimArray{D}} %DIF > 
\newcommand{\DistanceMatrix}{\BiDimArray{\Distance}} %DIF > 
\newcommand{\SubRate}{q} %DIF > 
\newcommand{\VarPhenotype}{V_{\Trait}} %DIF > 
\newcommand{\VarPhenotypeSpi}{V_{\Trait, \Spi}} %DIF > 
\newcommand{\VarGeneticSpi}{V_{\mathrm{A}, \Spi}} %DIF > 
\newcommand{\VarMutation}{V_{\mathrm{M}}} %DIF > 
\newcommand{\GenArchi}{\NbrLoci \Multiply \E \left[ a^2 \right]} %DIF > 
\newcommand{\RateMut}{\EstRateBrownian_{\mathrm{M}}} %DIF > 
\newcommand{\RateBrownian}{\EstRateBrownian} %DIF > 
\newcommand{\EstRateBrownian}{\hat{\sigma}^2} %DIF > 
\newcommand{\RootTrait}{\phi} %DIF > 
\newcommand{\EstRootTrait}{\hat{\RootTrait}} %DIF > 
\newcommand{\logit}{\text{logit}} %DIF > 
 %DIF > 
%DIF -------

\renewcommand{\baselinestretch}{1.3}
\renewcommand{\arraystretch}{1.2}
\linenumbers
\frenchspacing

\title{Phylogram instead of chronogram when assessing the neutral evolution of a trait}
\rhead{\scshape Phylogram over chronogram for neutral evolution.}

\author{
\large
\textbf{T. {Latrille}$^{1}$\orcidlink{0000-0002-9643-4668}, T. {Gaboriau}$^{1}$\orcidlink{0000-0001-7530-2204}, N. {Salamin}$^{1}$\orcidlink{0000-0002-3963-4954}}\\
\normalsize
$^{1}$Department of Computational Biology, University of Lausanne, Quartier Sorge, 1015 Lausanne, Switzerland\\
\texttt{\href{mailto:thibault.latrille@ens-lyon.org}{thibault.latrille@ens-lyon.org}} \\
}
%DIF PREAMBLE EXTENSION ADDED BY LATEXDIFF
%DIF UNDERLINE PREAMBLE %DIF PREAMBLE
\RequirePackage[normalem]{ulem} %DIF PREAMBLE
\RequirePackage{color}\definecolor{RED}{rgb}{1,0,0}\definecolor{BLUE}{rgb}{0,0,1} %DIF PREAMBLE
\providecommand{\DIFaddtex}[1]{{\protect\color{blue}\uwave{#1}}} %DIF PREAMBLE
\providecommand{\DIFdeltex}[1]{{\protect\color{red}\sout{#1}}} %DIF PREAMBLE
%DIF SAFE PREAMBLE %DIF PREAMBLE
\providecommand{\DIFaddbegin}{} %DIF PREAMBLE
\providecommand{\DIFaddend}{} %DIF PREAMBLE
\providecommand{\DIFdelbegin}{} %DIF PREAMBLE
\providecommand{\DIFdelend}{} %DIF PREAMBLE
\providecommand{\DIFmodbegin}{} %DIF PREAMBLE
\providecommand{\DIFmodend}{} %DIF PREAMBLE
%DIF FLOATSAFE PREAMBLE %DIF PREAMBLE
\providecommand{\DIFaddFL}[1]{\DIFadd{#1}} %DIF PREAMBLE
\providecommand{\DIFdelFL}[1]{\DIFdel{#1}} %DIF PREAMBLE
\providecommand{\DIFaddbeginFL}{} %DIF PREAMBLE
\providecommand{\DIFaddendFL}{} %DIF PREAMBLE
\providecommand{\DIFdelbeginFL}{} %DIF PREAMBLE
\providecommand{\DIFdelendFL}{} %DIF PREAMBLE
%DIF HYPERREF PREAMBLE %DIF PREAMBLE
\providecommand{\DIFadd}[1]{\texorpdfstring{\DIFaddtex{#1}}{#1}} %DIF PREAMBLE
\providecommand{\DIFdel}[1]{\texorpdfstring{\DIFdeltex{#1}}{}} %DIF PREAMBLE
\newcommand{\DIFscaledelfig}{0.5}
%DIF HIGHLIGHTGRAPHICS PREAMBLE %DIF PREAMBLE
\RequirePackage{settobox} %DIF PREAMBLE
\RequirePackage{letltxmacro} %DIF PREAMBLE
\newsavebox{\DIFdelgraphicsbox} %DIF PREAMBLE
\newlength{\DIFdelgraphicswidth} %DIF PREAMBLE
\newlength{\DIFdelgraphicsheight} %DIF PREAMBLE
% store original definition of \includegraphics %DIF PREAMBLE
\LetLtxMacro{\DIFOincludegraphics}{\includegraphics} %DIF PREAMBLE
\newcommand{\DIFaddincludegraphics}[2][]{{\color{blue}\fbox{\DIFOincludegraphics[#1]{#2}}}} %DIF PREAMBLE
\newcommand{\DIFdelincludegraphics}[2][]{% %DIF PREAMBLE
\sbox{\DIFdelgraphicsbox}{\DIFOincludegraphics[#1]{#2}}% %DIF PREAMBLE
\settoboxwidth{\DIFdelgraphicswidth}{\DIFdelgraphicsbox} %DIF PREAMBLE
\settoboxtotalheight{\DIFdelgraphicsheight}{\DIFdelgraphicsbox} %DIF PREAMBLE
\scalebox{\DIFscaledelfig}{% %DIF PREAMBLE
\parbox[b]{\DIFdelgraphicswidth}{\usebox{\DIFdelgraphicsbox}\\[-\baselineskip] \rule{\DIFdelgraphicswidth}{0em}}\llap{\resizebox{\DIFdelgraphicswidth}{\DIFdelgraphicsheight}{% %DIF PREAMBLE
\setlength{\unitlength}{\DIFdelgraphicswidth}% %DIF PREAMBLE
\begin{picture}(1,1)% %DIF PREAMBLE
\thicklines\linethickness{2pt} %DIF PREAMBLE
{\color[rgb]{1,0,0}\put(0,0){\framebox(1,1){}}}% %DIF PREAMBLE
{\color[rgb]{1,0,0}\put(0,0){\line( 1,1){1}}}% %DIF PREAMBLE
{\color[rgb]{1,0,0}\put(0,1){\line(1,-1){1}}}% %DIF PREAMBLE
\end{picture}% %DIF PREAMBLE
}\hspace*{3pt}}} %DIF PREAMBLE
} %DIF PREAMBLE
\LetLtxMacro{\DIFOaddbegin}{\DIFaddbegin} %DIF PREAMBLE
\LetLtxMacro{\DIFOaddend}{\DIFaddend} %DIF PREAMBLE
\LetLtxMacro{\DIFOdelbegin}{\DIFdelbegin} %DIF PREAMBLE
\LetLtxMacro{\DIFOdelend}{\DIFdelend} %DIF PREAMBLE
\DeclareRobustCommand{\DIFaddbegin}{\DIFOaddbegin \let\includegraphics\DIFaddincludegraphics} %DIF PREAMBLE
\DeclareRobustCommand{\DIFaddend}{\DIFOaddend \let\includegraphics\DIFOincludegraphics} %DIF PREAMBLE
\DeclareRobustCommand{\DIFdelbegin}{\DIFOdelbegin \let\includegraphics\DIFdelincludegraphics} %DIF PREAMBLE
\DeclareRobustCommand{\DIFdelend}{\DIFOaddend \let\includegraphics\DIFOincludegraphics} %DIF PREAMBLE
\LetLtxMacro{\DIFOaddbeginFL}{\DIFaddbeginFL} %DIF PREAMBLE
\LetLtxMacro{\DIFOaddendFL}{\DIFaddendFL} %DIF PREAMBLE
\LetLtxMacro{\DIFOdelbeginFL}{\DIFdelbeginFL} %DIF PREAMBLE
\LetLtxMacro{\DIFOdelendFL}{\DIFdelendFL} %DIF PREAMBLE
\DeclareRobustCommand{\DIFaddbeginFL}{\DIFOaddbeginFL \let\includegraphics\DIFaddincludegraphics} %DIF PREAMBLE
\DeclareRobustCommand{\DIFaddendFL}{\DIFOaddendFL \let\includegraphics\DIFOincludegraphics} %DIF PREAMBLE
\DeclareRobustCommand{\DIFdelbeginFL}{\DIFOdelbeginFL \let\includegraphics\DIFdelincludegraphics} %DIF PREAMBLE
\DeclareRobustCommand{\DIFdelendFL}{\DIFOaddendFL \let\includegraphics\DIFOincludegraphics} %DIF PREAMBLE
%DIF AMSMATHULEM PREAMBLE %DIF PREAMBLE
\makeatletter %DIF PREAMBLE
\let\sout@orig\sout %DIF PREAMBLE
\renewcommand{\sout}[1]{\ifmmode\text{\sout@orig{\ensuremath{#1}}}\else\sout@orig{#1}\fi} %DIF PREAMBLE
\makeatother %DIF PREAMBLE
%DIF COLORLISTINGS PREAMBLE %DIF PREAMBLE
\RequirePackage{listings} %DIF PREAMBLE
\RequirePackage{color} %DIF PREAMBLE
\lstdefinelanguage{DIFcode}{ %DIF PREAMBLE
%DIF DIFCODE_UNDERLINE %DIF PREAMBLE
  moredelim=[il][\color{red}\sout]{\%DIF\ <\ }, %DIF PREAMBLE
  moredelim=[il][\color{blue}\uwave]{\%DIF\ >\ } %DIF PREAMBLE
} %DIF PREAMBLE
\lstdefinestyle{DIFverbatimstyle}{ %DIF PREAMBLE
	language=DIFcode, %DIF PREAMBLE
	basicstyle=\ttfamily, %DIF PREAMBLE
	columns=fullflexible, %DIF PREAMBLE
	keepspaces=true %DIF PREAMBLE
} %DIF PREAMBLE
\lstnewenvironment{DIFverbatim}{\lstset{style=DIFverbatimstyle}}{} %DIF PREAMBLE
\lstnewenvironment{DIFverbatim*}{\lstset{style=DIFverbatimstyle,showspaces=true}}{} %DIF PREAMBLE
\lstset{extendedchars=\true,inputencoding=utf8}

%DIF END PREAMBLE EXTENSION ADDED BY LATEXDIFF

\begin{document}
%TC:ignore
\maketitle

% Abstract (≤ 300 words)
\begin{abstract}
At the species level, the evolution of traits is driven by a combination of selective and neutral forces. To disentangle these processes, different scenarios of evolution are modelled and compared. For quantitative traits under selection, the species is usually considered to track a trait optimum, and such an optimum can change along the branches of the species tree. On the other hand, neutral evolution is modelled with a trait changing randomly around the ancestral trait value along the different branches of the species tree. Regardless of the intricacy of modelling trait changes, the species tree is assumed to be known and obtained independently of the trait data analysed. Branch lengths are usually assumed to be in units proportional to time, and the tree is represented by a chronogram. The rationale is that time correlates with trait changes because of its direct connection with the number of generations that \DIFaddbegin \DIFadd{have }\DIFaddend occurred. However, since the generation time of species can also vary along the phylogenetic tree, we argue that their use introduces biases. In contrast, if the phylogenetic tree is represented by a phylogram with branch lengths in units of sequence divergence, this will account for the effect of changing generation time. In this study, we show using simulations that, for a trait evolving neutrally, the fit of a random evolution model has more support on a phylogram than on a chronogram. However, comparing models and testing different scenarios of selection using a phylogram leads to incorrect predictions. Given these results, we argue that we should use phylograms instead of chronograms when claiming that a trait is evolving under drift. Nevertheless, we support the \DIFdelbegin \DIFdel{fact that we should generally continue to use }\DIFdelend \DIFaddbegin \DIFadd{generally accepted use of }\DIFaddend chronograms to model selection acting on a quantitative trait.
\end{abstract}

\keywords{Comparative studies \and Morphological evolution \and Genetic drift \and Macroevolution \and Chronogram \and Phylogram}

\newpage
%TC:endignore
\section*{Lay summary}\label{sec:summary}
\DIFdelbegin \DIFdel{When we look at different species, we can see that some traits are different between them}\DIFdelend \DIFaddbegin \DIFadd{Many species, are characterised by trait differences between each other}\DIFaddend .
For example, the brain size of different primate and human lineages varies.
The question we ask is: \DIFdelbegin \DIFdel{does }\DIFdelend \DIFaddbegin \DIFadd{Does }\DIFaddend the time that passed explain the differences in their traits?
Or are genetic differences between species better at explaining the differences?
The standard in evolutionary biology is to use time to explain differences in traits.
However, if the trait is not under natural selection and instead is changing due to \DIFdelbegin \DIFdel{accumulating mutations}\DIFdelend \DIFaddbegin \DIFadd{the accumulation of neutral substitutions (because of genetic drift)}\DIFaddend , the genetic differences between species \DIFdelbegin \DIFdel{instead of time should best }\DIFdelend \DIFaddbegin \DIFadd{should better }\DIFaddend explain the differences in traits \DIFaddbegin \DIFadd{than the time since divergence}\DIFaddend .
As a result, we argue that using both genetic differences and species divergence time simultaneously provides a better understanding of how traits evolve across species.


\section*{Introduction}\label{sec:introduction}
% The model of neutral trait evolution, explained in the context of phylogenetic comparative methods.
By measuring changes in a phenotypic trait along a lineage, it is possible to infer the type of selection acting on it or whether this change is only due to random genetic drift.
At a larger scale, from the observed variations of traits across species, regimes of evolution are typically assessed using phylogenetic comparative methods, where traits are modelled as evolving along the branches of the species tree~\citep{felsenstein_phylogenies_1985, felsenstein_phylogenies_1988a, harmon_phylogenetic_2018}.
For example, to model neutral evolution, the mean trait value is said to follow a Brownian motion (BM), branching and evolving independently after each speciation event~\citep{felsenstein_phylogenies_1985, lynch_phenotypic_1986, hansen_translating_1996}.
In other words, for each branch of the tree, the value at the descendant node is normally distributed around the ancestral value, with a variance proportional to the branch length.
In this framework, reconstructing trait variation along the whole phylogeny as a BM \DIFaddbegin \DIFadd{process }\DIFaddend can thus constitute a null model of neutral trait evolution\DIFaddbegin \DIFadd{, although as we discuss later, BM can also approximate non-neutral evolutionary processes}\DIFaddend .

% Alternative models of trait evolution and their pitfalls.
Alternatively to a simple BM, adding a trend parameter in the BM is interpreted as a signature of directional selection at the phylogenetic scale\DIFaddbegin \DIFadd{, reflecting a consistent bias in trait evolution across lineages}\DIFaddend ~\citep{silvestro_early_2019}.
\DIFaddbegin \DIFadd{However, such trends can be difficult to distinguish from BM in analyses using only extant taxa without fossil calibration~}\citep{felsenstein_phylogenies_1988a}\DIFadd{.
}\DIFaddend More complex models to detect selection have been proposed, notably the Ornstein-Uhlenbeck (OU) processes, where trait variation is constrained around an optimum value, which is often interpreted as a signature of stabilizing selection~\citep{hansen_stabilizing_1997, butler_phylogenetic_2004, beaulieu_modeling_2012}.
Methodologically, this alternative model of evolution raises issues since an OU process might be statistically preferred over \DIFdelbegin \DIFdel{a }\DIFdelend BM due to sampling artifacts~\citep{silvestro_measurement_2015, cooper_cautionary_2016, price_detecting_2022}.
Adding biological realism, the OU process can be relaxed to allow for multiple optima along the phylogenetic tree, which is interpreted as a few abrupt changes in the environment along the phylogeny~\citep{ingram_surface_2013, uyeda_novel_2014, khabbazian_fast_2016, mitov_fast_2020, grabowski_cautionary_2023}.
Finally, the optimum can also be allowed to change continuously along the tree, which results in the trait itself reflecting the movement of the optimum along the lineages due to the constant adaptive evolution of the trait toward the optimum~\citep{hansen_translating_1996, hansen_three_2024}.
One special case of a continuously changing optimum is again the BM, where changes in the optimum are a consequence of the randomly changing environment~\citep{hansen_comparative_2008}, named fluctuating selection~\citep{holstad_evolvability_2024} or also diversifying selection due to the diversity of generated phenotypes at the clade level~\citep{latrille_detecting_2024}.
This modelling raises an issue in disentangling selection from neutral evolution: the \DIFdelbegin \DIFdel{best fit of BM can not be interpreted as the trait evolving neutrally, it can also be interpreted as a continuously moving optimum.
}\DIFdelend \DIFaddbegin \DIFadd{fact that BM is the best-supported model in a set does not necessarily imply that a trait is evolving neutrally.
Indeed, Brownian evolution can also result from non-neutral evolutionary processes, such as fluctuating natural selection.
More generally, due to the central limit theorem, the additive effects of many different neutral and non-neutral random factors might result in evolution that appears Brownian.
}\DIFaddend 

% Deriving the neutral expectation from standing variation
\DIFdelbegin \DIFdel{In order to }\DIFdelend \DIFaddbegin \DIFadd{To }\DIFaddend disentangle neutral evolution from selection, another approach is to contrast the observed rate of evolution to the neutral expectation~\citep{lande_genetic_1980}.
The neutral expectation can be obtained if the underlying genetic architecture of the trait is known and the trait \DIFaddbegin \DIFadd{is }\DIFaddend encoded by many loci of additive effects~\citep{barton_infinitesimal_2017, sella_thinking_2019}.
Additionally, we can also derive the expected changes in mean trait value along a lineage given the trait variance at the population scale~\DIFdelbegin %DIFDELCMD < \citep{turelli_heritable_1984, felsenstein_phylogenies_1988a}%%%
\DIFdelend \DIFaddbegin \citep{turelli_heritable_1984, turelli_phenotypic_1988, lynch_rate_1990, felsenstein_phylogenies_1988a}\DIFaddend .
When generalized to many species at the phylogenetic scale, contrasting between and within-species variations \DIFdelbegin \DIFdel{allow }\DIFdelend \DIFaddbegin \DIFadd{allows }\DIFaddend us to disentangle neutral evolution from selection\DIFdelbegin \DIFdel{~}%DIFDELCMD < \citep{latrille_detecting_2024}%%%
\DIFdel{.
}\DIFdelend \DIFaddbegin \DIFadd{.
Specifically, }\citet{latrille_detecting_2024} \DIFadd{showed that the ratio of between-species to within-species trait variation should scale linearly with nucleotide divergence for a neutral trait, providing a neutrality test that requires both types of variation data.
However, such within-species variation data are not always available, motivating the current study, which focuses solely on between-species variation and tree branch length units.
}\DIFaddend More generally, comparing the rate of evolution at different timescales allows inferring the regime of evolution~\citep{hansen_three_2024, holstad_evolvability_2024}.
However, trait variation at different timescales is not always available, and the genetic architecture of the trait is often unknown.
Instead, to disentangle neutral evolution from selection in the context of phylogenetic comparative methods, we hereby focus on the underlying tree and the unit of its branch lengths (Fig.~\ref{fig:methods}A).

%TC:ignore
\begin{figure*}[!htb]
    \centering
    \includegraphics[width=\textwidth] {figure-1}
    \caption{
        Panel~A: Across many species, the evolution of continuous traits can be modelled as a stochastic process evolving along the branches of a phylogenetic tree.
        The branches of such \DIFaddbeginFL \DIFaddFL{a }\DIFaddendFL phylogenetic tree can be measured in either time (chronogram) or number of substitutions (phylogram).
        Panel~B: Theoretically, changes of a trait under a moving optimum should be better predicted by a chronogram \DIFaddbeginFL \DIFaddFL{than by a phylogram, since the optimum is expected to change with time rather than with the number of generations}\DIFaddendFL .
        \DIFdelbeginFL \DIFdelFL{Instead}\DIFdelendFL \DIFaddbeginFL \DIFaddFL{Conversely}\DIFaddendFL , changes of a neutral trait should be better predicted by a phylogram \DIFaddbeginFL \DIFaddFL{than by a chronogram, since neutral trait changes are proportional to the number of generations (and thus nucleotide divergence), not time}\DIFaddendFL .
        \DIFaddbeginFL \DIFaddFL{Comparing the fit of a BM process on both trees can help disentangle selection from neutral evolution.
    }\DIFaddendFL }
    \label{fig:methods}
\end{figure*}
%TC:endignore

% The use of chronogram is widely accepted and de facto standard in phylogenetic comparative methods.
Typically, the tree is assumed to be known and obtained independently, and the \textit{de facto} standard in phylogenetic comparative methods is to use a chronogram, where the branch lengths are proportional to time~\citep{felsenstein_phylogenies_1985, harmon_phylogenetic_2018}.
Theoretically, for a neutrally evolving trait, trait changes depend directly on the number of generations~\citep{hansen_translating_1996}, which is proportional to time only if the average time between two consecutive generations, called generation time, is constant and equal for all species.
Since generation time depends on the time for an individual to reach sexual maturity, it can vary considerably between species, \DIFdelbegin \DIFdel{from a few months for coral reef pygmy gobies to 150 years for the Greenland shark~}%DIFDELCMD < \citep{demagalhaes_database_2009, nielsen_eye_2016}%%%
\DIFdelend \DIFaddbegin \DIFadd{even within a single clade.
For example, in mammals, generation time ranges from less than a year in mice to over 20 years in elephants~}\citep{demagalhaes_database_2009}\DIFadd{, representing variation that could substantially impact inferences of neutral trait evolution}\DIFaddend .
As a result, variations of generation time between species \DIFdelbegin \DIFdel{means }\DIFdelend \DIFaddbegin \DIFadd{suggest }\DIFaddend that modelling neutral evolution as \DIFdelbegin \DIFdel{a }\DIFdelend BM on a chronogram \DIFdelbegin \DIFdel{might induce }\DIFdelend \DIFaddbegin \DIFadd{induces }\DIFaddend biases~\citep{litsios_effects_2012}.

% Phylogram to model neutral trait evolution
Alternatively, on a phylogram, \DIFdelbegin \DIFdel{the }\DIFdelend branch lengths represent another quantity than time, which can be used as the backbone to model trait evolution and can potentially absorb changes in generation time.
For example, it is known that using phylograms in units of morphological distance can be used to improve ancestral trait reconstruction for a discrete character~\citep{wilson_chronogram_2022}.
From a genomic perspective, phylograms can also be in units of nucleotide divergence, that is, depicting the number of nucleotide substitutions occurring along a branch.
Because the number of neutral nucleotide substitutions accumulating along a branch is proportional to the number of generations~\citep{kimura_evolutionary_1968, kimura_neutral_1983}, nucleotide divergence would in this case absorb the effect of changing generation time (i.e. longer branches for lineages with shorter generation time).
Such a phylogram, if obtained from neutrally evolving genomic loci would be more appropriate to model a trait evolving neutrally~\citep{latrille_detecting_2024}.
\DIFaddbegin \DIFadd{More precisely, under the assumption of a constant genetic architecture, the covariance in mean trait value between a pair of species increases linearly with the number of shared generations, and thus with shared nucleotide divergence for neutral sequences (see supplementary material section~\ref{sec:theoretical-rationale} for the theoretical derivation).
}\DIFaddend Additionally, the \DIFdelbegin \DIFdel{rational }\DIFdelend \DIFaddbegin \DIFadd{rationale }\DIFaddend to use nucleotide divergence to absorb changes in generation time is also valid for changes in mutation rate (per loci and per generation) under the assumption of a constant genetic architecture of the trait and that changes in mutation rate impact the whole genome~\citep{latrille_detecting_2024}.
Empirically, such changes in mutation rate along a phylogenetic tree are also observed, although to a lesser extent than changes in generation time~\citep{bergeron_evolution_2023}.

% Phylogram in the case of selection
Under the alternative scenario involving selection, shorter generation time and higher mutation rate can also allow for the species to track the changing optimum faster and not lag behind~\citep{lande_natural_1976, hansen_comparative_2008}.
However, the timescale for this \DIFdelbegin \DIFdel{lag (hundreds or thousands of years}\DIFdelend \DIFaddbegin \DIFadd{adaptive lag (typically hundreds to thousands of generations}\DIFaddend ) is negligible compared to the timescale on which the optimum changes at the level of clades (millions of years) in a first approximation~\citep{hansen_three_2024}.
\DIFaddbegin \DIFadd{While this lag timescale is difficult to measure empirically, theoretical considerations suggest it should be much shorter than the macroevolutionary timescale.
}\DIFaddend Thus, at the phylogenetic scale, since shifts in the optimum fitness peak \DIFdelbegin \DIFdel{is }\DIFdelend \DIFaddbegin \DIFadd{are }\DIFaddend extrinsic to the species and \DIFdelbegin \DIFdel{is }\DIFdelend \DIFaddbegin \DIFadd{are }\DIFaddend dependent on the environment which varies with time, using a chronogram would be more accurate to model mean trait changes.
As a result, for a trait under selection, mean trait changes should be better explained by the time rather than the number of substitutions that occurred along a branch (Fig.~\ref{fig:methods}B).

% Here, we propose a method to evaluate the soundness of studying trait evolution on a phylogram or a chronogram.
In this study, we test this hypothesis, namely \DIFdelbegin \DIFdel{if }\DIFdelend \DIFaddbegin \DIFadd{whether }\DIFaddend using a phylogram (in units of nucleotide divergence for neutral loci) instead of a chronogram (in units proportional to time) can allow to assess the \DIFdelbegin \DIFdel{regime of evolution of }\DIFdelend \DIFaddbegin \DIFadd{evolutionary regime of }\DIFaddend a trait, by evaluating the fit of \DIFdelbegin \DIFdel{a }\DIFdelend BM on both trees.
More precisely, we seek to test \DIFdelbegin \DIFdel{if }\DIFdelend \DIFaddbegin \DIFadd{whether }\DIFaddend the changes of a neutral trait are better predicted by a phylogram.
Conversely, we also seek to test \DIFdelbegin \DIFdel{if }\DIFdelend \DIFaddbegin \DIFadd{whether }\DIFaddend changes of a trait under a moving optimum are better predicted by a chronogram.
Using genomic information\DIFaddbegin \DIFadd{, }\DIFaddend we seek to provide additional approaches to model trait evolution, particularly to disentangle selection and neutral evolution.

\section*{Methods}
To assess whether a phylogram is better suited to model neutral evolution\DIFaddbegin \DIFadd{, }\DIFaddend we first performed simulations of a trait evolving along a phylogenetic tree under different regimes of selection (neutral, moving optimum, multiple optima).
Second, we fitted a \DIFdelbegin \DIFdel{single rate }\DIFdelend \DIFaddbegin \DIFadd{single-rate }\DIFaddend Brownian motion (BM) to the simulated data, and we compared the fit of the BM on a phylogram versus a chronogram.
Third, based on the \DIFdelbegin \DIFdel{single rate }\DIFdelend \DIFaddbegin \DIFadd{single-rate }\DIFaddend BM, we derived a model to estimate the support of a phylogram over a chronogram that we apply to the simulated dataset.
Fourth, we fitted several models of trait evolution to the simulated datasets: a multi-rate BM, an Ornstein-Uhlenbeck (OU) process, and a relaxed OU with multiple optima, tested on a phylogram versus a chronogram.
Finally, we gathered an empirical dataset of body and brain masses from mammals, including both a chronogram and a phylogram (in units of nucleotide divergence from neutral loci), on which we tested the support of a phylogram for the \DIFdelbegin \DIFdel{single rate BM.
}\DIFdelend \DIFaddbegin \DIFadd{single-rate BM.
To assess which tree type (phylogram versus chronogram) better supports neutral evolution, we evaluate both the statistical fit of the BM and the accuracy of ancestral trait reconstruction under each tree.
}\DIFaddend 

\subsection*{Simulations along a phylogenetic tree}
We performed simulations under different selective regimes (neutral, moving optimum, multiple optima).
Simulations were individual-based and followed a Wright-Fisher model with mutation, selection and drift for a diploid population including speciation along a predefined ultrametric phylogenetic tree.
We used the same simulation framework as in \citet{latrille_detecting_2024}, with parameters detailed in the supplementary material (section~\ref{sec:simulator}).
The parameters of simulations were chosen to mimic an empirical dataset of mammals.
In summary, the trait was encoded by $\NbrLoci$ independent loci, with each locus contributing additively, and mutations were drawn from a Poisson distribution at each generation.
Parents were selected for reproduction according to their phenotypic value, with a probability proportional to their fitness.
Flattening the fitness landscape resulted in neutral evolution, which meant that each individual had the same probability of being sampled at each generation regardless of its trait value.
Alternatively, for a trait under selection around a moving optimum, we modelled stabilizing selection acting on the trait, with the optimum value changing as a geometric Brownian motion (BM) along the phylogenetic tree~\citep{hansen_stabilizing_1997, hansen_translating_1996}.
Finally, for a trait under selection around multiple optima~\citep{uyeda_novel_2014}, we draw a Bernoulli variable to test whether there is a switch of optimum along a branch; if a switch occurs, the sliding of the optimum is drawn from a reflected exponential distribution (symmetric positive and negative values).

At each internal node of the tree (i.e. speciation event), the population is split into two daughter populations running independently on each of the two branches, and the process is repeated until the tips of the tree are reached.
As a control, we performed simulations with constant mutation rates, generation times and effective population sizes ($\Ne$).
Alternatively, we performed simulations with fluctuating mutation rates, generation times and $\Ne$, where we used \DIFdelbegin \DIFdel{a }\DIFdelend BM to model the long-term changes along the phylogenetic tree, and we overlaid short-term changes in $\Ne$ (see supplementary material section~\ref{sec:simulator}).

The phylogram is obtained from the same simulation setting on a set of 30,000 independent neutral loci, and the nucleotide divergence is computed as the number of substitutions along each branch.
The chronogram is \DIFaddbegin \DIFadd{the generating ultrametric time tree used in the simulation.
For the empirical datasets, the chronogram is }\DIFaddend derived from the phylogram by fitting a relaxed molecular clock model (correlated rate model with penalized likelihood, default value) to the phylogram with the \textit{R} package \textit{ape}~\citep{paradis_ape_2004}.

\subsection*{Brownian motion (BM)}
On the simulated dataset, we fitted a \DIFdelbegin \DIFdel{single rate }\DIFdelend \DIFaddbegin \DIFadd{single-rate }\DIFaddend BM using either a phylogram or a chronogram as the underlying tree.
We used \textit{RevBayes}~\citep{hohna_revbayes_2016} to fit a Brownian motion (BM) on continuous traits evolving along the branches of a phylogenetic tree~\citep{felsenstein_phylogenies_1985, felsenstein_phylogenies_1988a}.
The data consist of the mean trait value for each extant species, and the tree topology is fixed.
Along each branch, the value of the trait is drawn from a normal distribution with mean equal to the parent node value and variance equal to the branch length~\citep{felsenstein_phylogenies_1985}.
Formally, the BM process is described by the equation:
% Ornstein-Uhlenbeck model equations
\begin{align}
    \der \Trait_t & = \sigma \Multiply \der W_t,  \label{eq:bm-process}
\end{align}
where $\Trait_t$ is the trait value at time $t$, $\sigma$ is the rate of the BM, and $\der W_t$ is a Wiener process (i.e. a standard Brownian motion).
We used a log-uniform prior for $\sigma$.
Using Markov chain Monte Carlo (MCMC), we reconstructed the ancestral trait value at each node of the tree as the posterior mean estimate (burn-in of 1000 gen., running of 10,000 gen., 2 chains), and we compared the accuracy of the reconstruction using either a phylogram or a chronogram as the underlying tree.
Both trees are scaled such that the sum of all the branch lengths is equal to one in each case.
Importantly, the data and the tree topology are the same in both analyses; only the branch lengths are different.
In this simulation setting, for traits simulated under neutral evolution, we expect more accurate trait reconstruction on a phylogram, and conversely, for traits under selection (moving optimum), more accurate trait reconstruction on a chronogram.

\subsection*{BM with a switch}
To test the support of a phylogram over a chronogram, we implemented a model based on a \DIFdelbegin \DIFdel{single rate }\DIFdelend \DIFaddbegin \DIFadd{single-rate }\DIFaddend BM.
The model was implemented in \textit{RevBayes} and contains a switch variable, denoted as $\PhyloSupport$ that allows the model to switch the branch lengths in units of time to units of substitutions (Fig.~\ref{fig:results-brownian}A).
Formally, \DIFdelbegin \DIFdel{$\PhyloSupport \in {0, 1}$ }\DIFdelend \DIFaddbegin \DIFadd{$\PhyloSupport$ }\DIFaddend is a Bernoulli random variable (\DIFdelbegin \DIFdel{$\PhyloSupport \in {0, 1}$}\DIFdelend \DIFaddbegin \DIFadd{$\PhyloSupport \in \{0, 1\}$}\DIFaddend ) with a prior probability of $0.5$.
If $\PhyloSupport$ is $0$, the branch lengths are those of the chronogram, and if $\PhyloSupport$ is $1$, the branch lengths are those of the phylogram.
The tree topology is fixed, meaning both trees have the same branching structure, but because branch lengths are different and not necessarily on the same scale, both trees are re-scaled by dividing branch lengths by the total tree length.
As in the previous section, \DIFdelbegin \DIFdel{a }\DIFdelend BM is fitted to the data by modelling trait changes as normal distributions (\textit{dnNormal} in \textit{RevBayes}) running along the branches of the tree, with variance proportional to the branch length and the squared value of the Brownian rate parameter ($\sigma$, log-uniform prior).
Alternatively for large trees or for faster computation, the likelihood can also be estimated by the REML method (\textit{dnPhyloBrownianREML} in \textit{RevBayes}).
Altogether, the input data is the mean trait value for extant species and both the chronogram and phylogram with the same tree topology.
The posterior mean of $\PhyloSupport$ (burn-in of 1,000 gen., running of 10,000 gen., 2 chains) \DIFdelbegin \DIFdel{is the probability that the phylogram is favoured over the chronogram}\DIFdelend \DIFaddbegin \DIFadd{indicates support for the phylogram; we consider the phylogram to be favoured when $\PhyloSupport > 0.5$}\DIFaddend .
We expect $\PhyloSupport = 1$ for a trait under neutral evolution, and $\PhyloSupport = 0$ for a trait under selection (moving optimum).

\DIFaddbegin \DIFadd{From a frequentist perspective, the support of a phylogram over a chronogram can also be assessed by comparing the maximum likelihood of the data given each tree, and computing the Akaike Information Criterion (AIC) weights of each model (see supplementary material section~~\ref{sec:maximum-likelihood-estimation}).
}

\DIFaddend \subsection*{Alternative models of evolution}
In addition to a \DIFdelbegin \DIFdel{single rate }\DIFdelend \DIFaddbegin \DIFadd{single-rate }\DIFaddend BM, we also fitted \DIFdelbegin \DIFdel{a }\DIFdelend BM with multiple rate parameters~\citep{eastman_novel_2011}, using either a phylogram or a chronogram as the underlying tree to the simulated dataset.
The likelihood of the data is estimated by the REML method (\textit{dnPhyloBrownianREML} in \textit{RevBayes}), and the number of rate shifts is estimated by reversible-jump MCMC (\textit{dnReversibleJumpMixture} in \textit{RevBayes}).
For each branch, we draw a rate-multiplier either equal to $1$ (no rate shift), or drawn from a log-normal distribution with a median of 1, and a standard deviation such that rate shifts range over about one order of magnitude.
We set our prior to the expected number of rate shifts of $1$ across the whole tree.
The number of rate shifts as well as the variance of rate parameters are estimated as the posterior mean (burn-in of 1,000 gen., running of 50,000 gen., 2 chains).

% Simple OU model versus a brownian process.
Also, we compared the fit of an Ornstein-Uhlenbeck (OU) process (eq.~\ref{eq:ou-process}) to \DIFdelbegin \DIFdel{a }\DIFdelend BM~\citep{hansen_stabilizing_1997,butler_phylogenetic_2004}, using either a phylogram or a chronogram as the underlying tree.
% Ornstein-Uhlenbeck model equations
The OU process is described by the equation:
\begin{align}
    \der \Trait_t & = -\alpha \left( \Trait_t - \theta \right) \der t + \sigma \Multiply \der W_t, \label{eq:ou-process}
\end{align}
where the parameter $\alpha$ in eq.~\ref{eq:ou-process} is the strength of the pull towards the optimum, $\theta$ is the optimum value and the other parameters are defined as in eq~\ref{eq:bm-process}.
The likelihood of the data is estimated by the REML method (\textit{dnPhyloOrnsteinUhlenbeckREML} in \textit{RevBayes}), and because the OU process has more parameters than the BM, we used a reversible-jump MCMC switch between the two models ($0$ for BM and $1$ for OU, \textit{dnReversibleJumpMixture} in \textit{RevBayes}).
% alpha ~ dnReversibleJumpMixture(0.0, dnExponential( abs(root_age / 2.0 / ln(2.0)) ), 0.5)
The prior for $\theta$ is uniform between the minimum and maximum trait value for extant species.
The prior for $\alpha$ is a reversible-jump mixture distribution: a value of $0$ or drawn from an exponential distribution with mean equal to half the root age divided by $\ln (2)$, meaning that we expect a phylogenetic \DIFdelbegin \DIFdel{half life }\DIFdelend \DIFaddbegin \DIFadd{half-life }\DIFaddend of half the \DIFdelbegin \DIFdel{tree }\DIFdelend \DIFaddbegin \DIFadd{crown }\DIFaddend age.
Thus, when $\alpha = 0$ the OU process (eq.~\ref{eq:ou-process}) is equivalent to \DIFdelbegin \DIFdel{a }\DIFdelend BM (eq.~\ref{eq:bm-process}).
The \DIFdelbegin \DIFdel{support for the OU model is }\DIFdelend \DIFaddbegin \DIFadd{half-life, $t_{1/2}$, is defined as $\ln(2) / \alpha$ is also }\DIFaddend estimated as the posterior mean \DIFdelbegin \DIFdel{of the switch variable ($p_{\text{OU}}$), or equivalently that of $\alpha \neq 0$ (}\DIFdelend \DIFaddbegin \DIFadd{estimate of finite values (}\DIFaddend burn-in of 1,000 gen., running of \DIFdelbegin \DIFdel{10}\DIFdelend \DIFaddbegin \DIFadd{5}\DIFaddend ,000 gen., 2 chains)\DIFdelbegin \DIFdel{.
}\DIFdelend \DIFaddbegin \DIFadd{, representing the time needed for the expected trait value to move half-way toward the optimum~}\citep{hansen_stabilizing_1997, grabowski_cautionary_2023}\DIFadd{,
}\DIFaddend 

% Relaxed OU model, how many optimums?
Finally, we fitted a relaxed OU with multiple optima~\citep{uyeda_novel_2014}, using either a phylogram or a chronogram as the underlying tree.
The likelihood of the data given the tree is estimated by the REML method (\textit{dnPhyloOrnsteinUhlenbeckREML} in \textit{RevBayes}), and the number of optimums is also estimated by the reversible-jump MCMC method (\textit{dnReversibleJumpMixture} in \textit{RevBayes}) as the posterior mean estimate (burn-in of 1,000 gen., running of 5,000 gen., 2 chains).
For each branch, we draw a shift in optimum either equal to $0$ (no shift), or drawn from a uniform distribution centred on $0$ with breadth spanning the whole range of observed trait value.
We set our prior to the expected number of optimum shifts of $1$ across the whole tree.

\subsection*{Empirical dataset}
We analysed a dataset of body and brain masses from mammals.
The log-transformed values of body and brain masses were taken from \citet{tsuboi_breakdown_2018}.
We removed individuals not marked as adults and split the data into males and females due to sexual dimorphism in body and brain masses.
We also extracted brain and body masses from the \textit{COMBINE} dataset~\citep{soria_combine_2021}.
The mammalian genomic data are gathered from the Zoonomia project~\citep{genereux_comparative_2020}.
More specifically, the phylogram in units of nucleotide divergence is estimated on a set of neutral markers in \citet{foley_genomic_2023}.
The chronogram is derived from the phylogram by fitting a relaxed molecular clock model (correlated rate model with penalized likelihood, default value) to the phylogram with the \textit{R} package \textit{ape}~\citep{paradis_ape_2004}.

%TC:ignore
\begin{figure*}[!htb]
    \centering
    \includegraphics[width=\textwidth] {figure-2}
    \caption{
        \textbf{Testing the support for a phylogram over a chronogram}
        Panel~A:
        Brownian motion (BM) is used to model trait changes along a phylogeny, given a mean trait value for extant species.
        The model contains a switch variable, denoted $\PhyloSupport$, that switches the branch lengths in units of time ($\PhyloSupport=0$, chronogram) to units of nucleotide divergence for neutral sites ($\PhyloSupport=1$, phylogram).
        The posterior mean estimate of $\PhyloSupport$ indicates whether \DIFdelbeginFL \DIFdelFL{a }\DIFdelendFL BM supports \DIFdelbeginFL \DIFdelFL{better }\DIFdelendFL the phylogram \DIFdelbeginFL \DIFdelFL{over }\DIFdelendFL \DIFaddbeginFL \DIFaddFL{better than }\DIFaddendFL the chronogram ($0 \leq \PhyloSupport \leq 1$).
        Panel~B: Wright-Fisher simulations of trait evolution along a phylogeny.
        At each node of the tree a speciation event creates two descendant species evolving independently.
        Neutral evolution is modelled as a constant fitness regardless of the phenotype of individuals (yellow).
        Selection is modelled by stabilizing selection around \DIFdelbeginFL \DIFdelFL{a }\DIFdelendFL \DIFaddbeginFL \DIFaddFL{an }\DIFaddendFL optimum value, which is changing along the phylogenetic tree: as a moving optimum (blue) or as multiple optima (red).
        The output of the simulation is: 1) the mean trait value for each extant species, 2) the phylogram as obtained from independent neutral markers and 3) the chronogram derived from the phylogram assuming a relaxed molecular clock.
        The posterior mean estimate of $\PhyloSupport$ is inferred from the output of the simulation.
        Panels C \& D: Violin plot of the posterior mean of $\PhyloSupport$ across 100 replicates for the different simulated regimes of evolution with mean $\PhyloSupport$ values at the center ($\MeanPhyloSupport$).
        Horizontal lines inside the violins are the estimated $\PhyloSupport$ of each replicate \DIFdelbeginFL \DIFdelFL{simulations}\DIFdelendFL \DIFaddbeginFL \DIFaddFL{simulation}\DIFaddendFL .
        Results of simulations with constant generation time (Panel~C) and with fluctuating generation time, mutation rate and effective population size (Panel~D).
    }
    \label{fig:results-brownian}
\end{figure*}
%TC:endignore

\section*{Results}

\subsection*{Support for a phylogram over a chronogram}
% Simulations of a trait evolving neutrally versus under a moving optimum.
We tested the hypothesis that a phylogram is better suited to model neutral evolution than a chronogram.
We simulated traits evolving on a phylogenetic tree under different regimes of selection: 1) neutrally evolving, 2) under a moving optimum, and 3) under multiple optima (see Methods).
We first assessed the accuracy of ancestral trait reconstruction on a phylogram versus a chronogram, under the assumption that changes follow \DIFdelbegin \DIFdel{a }\DIFdelend BM.
For a neutral trait, ancestral trait reconstruction was more accurate on a phylogram (Fig.~\ref{fig:results-distance}A, $r^2=0.95$) than on a chronogram (Fig.~\ref{fig:results-distance}B, \DIFdelbegin \DIFdel{$r^2=0.85$}\DIFdelend \DIFaddbegin \DIFadd{$r^2=0.83$; Fisher's $z$-test, $p_{\text{value}}=9.0\times 10^{-11}$}\DIFaddend ).
In contrast, for a trait evolving under a moving optimum, ancestral trait reconstruction was less accurate on a phylogram (Fig.~\ref{fig:results-distance}C, $r^2=0.79$) than on a chronogram (Fig.~\ref{fig:results-distance}D, $r^2=0.96$\DIFaddbegin \DIFadd{; Fisher's $z$-test, $p_{\text{value}}=0$}\DIFaddend ).

We then tested for the BM support of phylogram over a chronogram, and we developed a statistic denoted as $\PhyloSupport$ (Fig.~\ref{fig:results-brownian}A), which ranges from 0 (full support for the chronogram) to 1 (full support for the phylogram).
For simulations with a constant generation time, the support is similar for both a phylogram and a chronogram regardless of the regime of evolution (Fig.~\ref{fig:results-brownian}B-C), which is expected since the branch lengths are equivalent in both trees.
Next, we simulated changing generation time, mutation rate (per generation) and effective population size ($\Ne$) along the phylogenetic tree, mimicking a mammalian range of changes.
In this case, for a simulated neutral trait, the BM was better fitted on a phylogram than on a chronogram, with an average support for phylogram of \DIFdelbegin \DIFdel{$\MeanPhyloSupport = 0.94$ }\DIFdelend \DIFaddbegin \DIFadd{$\MeanPhyloSupport = 0.95$ }\DIFaddend across $100$ replicate simulations, and \DIFdelbegin \DIFdel{$80$ }\DIFdelend \DIFaddbegin \DIFadd{$82$ }\DIFaddend of the $100$ replicates above the $0.95$ threshold (Fig.~\ref{fig:results-brownian}D, yellow).
Conversely, for a simulated trait under selection, the fit of \DIFdelbegin \DIFdel{a }\DIFdelend BM on a chronogram was better than on a phylogram with an average support for phylogram of \DIFdelbegin \DIFdel{$\MeanPhyloSupport = 0.17$ }\DIFdelend \DIFaddbegin \DIFadd{$\MeanPhyloSupport = 0.12$ }\DIFaddend for multiple optima (\DIFdelbegin \DIFdel{$69$ }\DIFdelend \DIFaddbegin \DIFadd{$67$ }\DIFaddend out of $100$ replicates below the $0.05$ threshold, Fig.~\ref{fig:results-brownian}D, blue) and \DIFdelbegin \DIFdel{$\MeanPhyloSupport = 0.14$ }\DIFdelend \DIFaddbegin \DIFadd{$\MeanPhyloSupport = 0.24$ }\DIFaddend for a moving optimum (\DIFdelbegin \DIFdel{$67$ }\DIFdelend \DIFaddbegin \DIFadd{$62$ }\DIFaddend of the $100$ below the $0.05$ threshold, Fig.~\ref{fig:results-brownian}D, red).
\DIFaddbegin \DIFadd{Finally, from a statistical perspective, the equivalent of $\PhyloSupport$ using the maximum likelihood approach (AIC weights, see supplementary material section~\ref{sec:maximum-likelihood-estimation}) is highly correlated with our Bayesian estimate ($r^2=0.99$, Fig.~\ref{fig:aicweights}).
}\DIFaddend 

On the empirical mammalian dataset from \citet{tsuboi_breakdown_2018}, for body mass the support for the phylogram is $\PhyloSupport = 0.0$, when sex is not taken into account.
When splitting the dataset into males and females, the support for the phylogram is $\PhyloSupport_{\text{\Female}} = 0.85$ and $\PhyloSupport_{\text{\Male}} = 0.70$.
For brain mass, the support for the phylogram is $\PhyloSupport = 0.0035$ on the mixed dataset, with $\PhyloSupport_{\text{\Female}} = 0.56$ and $\PhyloSupport_{\text{\Male}} = 0.51$.
On the \textit{COMBINE} dataset, which does not distinguish between males and females, the support for the phylogram is $\PhyloSupport = 0.0$ for body mass and $\PhyloSupport = 0.0015$ for brain mass (Table~\ref{table:results-empirical}).
We also computed $\PhyloSupport$ using an approximate likelihood computation (REML), a method that allow for faster computation and scale better for larger trees, and found similar results (Table~\ref{table:results-empirical}).

\subsection*{Fitting alternative models of evolution}
% Mis-classification of trait evolution on a phylogram.
Besides testing the fit of \DIFdelbegin \DIFdel{a }\DIFdelend BM with a \DIFdelbegin \DIFdel{single rate}\DIFdelend \DIFaddbegin \DIFadd{single-rate}\DIFaddend , we also fitted a multi-rate BM (see Methods) \DIFaddbegin \DIFadd{to simulated datasets with changing generation time, mutation rate and effective population size ($\Ne$)}\DIFaddend .
When fitting a multi-rate BM, the estimated variance of rate parameters ($v$) was lower on a phylogram than on a chronogram for a trait evolving neutrally (Fig.~\ref{fig:results-alternative}A, yellow violins, Wilcoxon paired test with \DIFdelbegin \DIFdel{$p_{\text{value}}=8.8\times 10^{-14}$}\DIFdelend \DIFaddbegin \DIFadd{$p_{\text{value}}=4.5\times 10^{-13}$}\DIFaddend ), and the number of rate shifts ($n$) was reflecting the prior on the phylogram while not on the chronogram, showing a bias (Fig.~\ref{fig:results-alternative}B, yellow violins).
Conversely, for a trait under a moving optimum, $v$ was higher on a phylogram than on a chronogram (Fig.~\ref{fig:results-alternative}A, blue violins, Wilcoxon paired test with  \DIFdelbegin \DIFdel{$p_{\text{value}}=8.9\times 10^{-15}$}\DIFdelend \DIFaddbegin \DIFadd{$p_{\text{value}}=1.3\times 10^{-14}$}\DIFaddend ), and $n$ was reflecting the prior on the chronogram while not on the phylogram, showing a bias (Fig.~\ref{fig:results-alternative}B, blue violins).
Altogether, fitting \DIFdelbegin \DIFdel{a }\DIFdelend BM on a phylogram is more accurate for a trait evolving neutrally, but results in less accurate estimates for a trait under selection.

% Mainly testing BM versus OU
Moreover, we tested alternatives to the BM with OU models: with a single optimum and with multiple optima (see Methods).
\DIFdelbegin \DIFdel{First, }\DIFdelend \DIFaddbegin \DIFadd{Following recommendations by }\citet{grabowski_cautionary_2023}\DIFadd{, we report the phylogenetic half-life ($t_{1/2} = \ln(2)/\alpha$), which represents the time }\DIFaddend for a trait \DIFdelbegin \DIFdel{under a moving optimum, the estimated support for the OU process ($p_{\text{OU}}$) was higher }\DIFdelend \DIFaddbegin \DIFadd{to evolve halfway toward a new optimum.
The higher the half-life, or in other words, the smaller the $\alpha$, the more the OU process becomes effectively indistinguishable from BM.cd m
For a trait evolving neutrally, the estimated half-life was longer }\DIFaddend on a phylogram than on a chronogram (Fig.~\ref{fig:results-alternative}C, \DIFdelbegin \DIFdel{blue }\DIFdelend \DIFaddbegin \DIFadd{yellow }\DIFaddend violins, Wilcoxon paired test with \DIFdelbegin \DIFdel{$p_{\text{value}}=3.9\times 10^{-18}$), an expected result since the BM with a switch supported the chronogram.
For a trait evolving neutrally, $p_{\text{OU}}$ was still higher on a phylogram }\DIFdelend \DIFaddbegin \DIFadd{$p_{\text{value}}=4.9\times 10^{-5}$), confirming BM-like dynamics on phylograms as expected for neutral evolution.
Conversely, for a trait under a moving optimum, the half-life was shorter on a chronogram }\DIFaddend than on a \DIFdelbegin \DIFdel{chronogram }\DIFdelend \DIFaddbegin \DIFadd{phylogram }\DIFaddend (Fig.~\ref{fig:results-alternative}C, \DIFdelbegin \DIFdel{yellow }\DIFdelend \DIFaddbegin \DIFadd{blue }\DIFaddend violins, Wilcoxon paired test with \DIFdelbegin \DIFdel{$p_{\text{value}}=8.0\times 10^{-8}$), an unexpected result since the BM with a switch supported the phylogram, thus showing a high rate of mis-classification on a phylogram}\DIFdelend \DIFaddbegin \DIFadd{$p_{\text{value}}=1.5\times 10^{-12}$), indicating faster adaptation on the appropriate tree}\DIFaddend .
When fitting a multi-OU process on a trait evolving neutrally, the number of optimum shifts ($m$) was reflecting the prior on the phylogram, but not on the chronogram (Fig.~\ref{fig:results-alternative}\DIFdelbegin \DIFdel{B}\DIFdelend \DIFaddbegin \DIFadd{D}\DIFaddend , yellow violins).
Conversely, for a trait under a moving optimum, ($m$) was reflecting the prior on the chronogram but not on the phylogram (Fig.~\ref{fig:results-alternative}\DIFdelbegin \DIFdel{A}\DIFdelend \DIFaddbegin \DIFadd{D}\DIFaddend , blue violins).
\DIFdelbegin \DIFdel{In summary, fitting an OU process on a phylogram results in biases not only for a trait under selection but also for a trait under neutral evolution.
}\DIFdelend 


%TC:ignore
\begin{figure*}[!htb]
    \centering
    \includegraphics[width=\textwidth] {figure-3}
    \caption{
        \textbf{Mis-classification of trait evolution on a phylogram.}
        Violin plot of posterior parameter estimates for different models of trait evolution and different simulated regimes of selection.
        Horizontal lines inside the violin show the results of each replicate simulation.
        Simulations of 100 replicates per regime: trait evolving under a neutral regime (yellow) and under a moving optimum (blue).
        Wilcoxon rank-test are performed between the paired estimates on the phylogram and the chronogram.
        Panel~A \& B: Fit of a relaxed Brownian motion (BM) with multiple rate parameters, using either a phylogram or a chronogram.
        Posterior estimate for the variance in rate parameters (Panel~A) and number of rate changes (Panel~B).
        Panel~C: \DIFdelbeginFL \DIFdelFL{Relative fit of }\DIFdelendFL \DIFaddbeginFL \DIFaddFL{Estimated phylogenetic half-life ($t_{1/2} = \ln(2)/\alpha$) from }\DIFaddendFL an Ornstein-Uhlenbeck (OU) process\DIFdelbeginFL \DIFdelFL{compared to BM}\DIFdelendFL , using either a phylogram or a chronogram.
        \DIFdelbeginFL \DIFdelFL{Posterior estimates }\DIFdelendFL \DIFaddbeginFL \DIFaddFL{The half-life represents the time }\DIFaddendFL for \DIFaddbeginFL \DIFaddFL{a trait to evolve halfway toward }\DIFaddendFL the \DIFdelbeginFL \DIFdelFL{support of the OU process over the BM (reversible-jump)}\DIFdelendFL \DIFaddbeginFL \DIFaddFL{optimum; longer half-life indicates BM-like dynamics}\DIFaddendFL .
        Panel~D: Fit of a relaxed OU with multiple optima, using either a phylogram or a chronogram.
        Posterior estimate for the number of optimum changes.
    }
    \label{fig:results-alternative}
\end{figure*}
%TC:endignore

\section*{Discussion}
% Context and rationale
In phylogenetic comparative methods, the evolution of continuous traits is typically modelled as a stochastic process \DIFdelbegin \DIFdel{running }\DIFdelend \DIFaddbegin \DIFadd{evolving }\DIFaddend along the branches of a phylogenetic tree.
The null model of neutral trait evolution, genetic drift, is thought of as a Brownian motion (BM) running on the tree~\citep{lynch_phenotypic_1986, felsenstein_phylogenies_1988a}, and deviations from this model are interpreted as selection acting on the trait~\citep{butler_phylogenetic_2004}.
As a result, a trait on which the BM is favoured over alternative models are sometimes interpreted as a trait evolving neutrally~\citep{khaitovich_evolution_2006, catalan_drift_2019}.
However, the BM is not only the null model of neutral evolution, but can also model selection, with a trait tracking a moving optimum, where the optimum itself is following \DIFdelbegin \DIFdel{a }\DIFdelend BM~\citep{hansen_translating_1996, latrille_detecting_2024}.
\DIFaddbegin \DIFadd{More fundamentally, the central limit theorem implies that trait evolution can appear Brownian whenever the trait is influenced by a large number of independent random factors acting additively~}\citep{felsenstein_phylogenies_1985}\DIFadd{.
These factors may include fluctuating genetic architecture, mutational pressures, changing population size, randomly changing environment, genetic drift, and natural selection.
Some of these influences may vary in a manner more directly connected to generation time, while others may track real time, making it difficult to predict a priori which tree type should provide the better fit.
Therefore, a good fit of BM alone cannot distinguish between neutral evolution and selection, and the comparison between phylogram and chronogram provides a crucial additional diagnostic.
}

\DIFaddend Distinguishing these scenarios is not trivial, but we here showed that the underlying backbone tree to model trait evolution can help disentangle neutral evolution from selection.
The standard in phylogenetic comparative methods is a chronogram, where the branch lengths are measured in time~\citep{felsenstein_phylogenies_1985, harmon_phylogenetic_2018}, while phylograms, where the branch lengths are measured in number of substitutions instead, have been generally overlooked.
For a simulated trait under selection we assessed that the BM has indeed a better fit on a chronogram than on a phylogram.
However, we argue that to model a trait evolving neutrally, the backbone tree should not be a chronogram but instead a phylogram.
Phylograms better represent the number of generations that occurred, and in turn should explain best the differences in traits between species~\citep{latrille_detecting_2024}
In practice, we show that for simulations of a trait evolving neutrally, the fit of \DIFdelbegin \DIFdel{a }\DIFdelend BM supports a phylogram rather than a chronogram.
Additionally, for a neutral trait, both ancestral state reconstruction and estimation for the number of rate changes for a multi-rate BM was more accurate on a phylogram.
Altogether, the combined use of a phylogram and a chronogram can help support the hypothesis of neutral evolution, or to rule it out.
% Thus, we argue that claims for a trait evolving under drift should be tested more thoroughly by comparing the fit of the BM under both a phylogram and a chronogram as the underlying tree.

% When does it apply and why is it important
We applied this method to different empirical datasets of body and brain masses in mammals, and we showed that the chronogram was favoured over the phylogram, ruling out the neutral model of evolution for these traits~\citep{latrille_detecting_2024}.
The mammalian dataset is particularly relevant, since changes in generation time are expected to occur along the phylogenetic tree, from 6 months for mice to 50 years for bowhead whales~\citep{demagalhaes_database_2009, nielsen_eye_2016}.
\DIFdelbegin \DIFdel{But more }\DIFdelend \DIFaddbegin \DIFadd{When analysing each sex separately, the dataset with annotated sex available is smaller (column Taxa in Table~\ref{table:results-empirical}).
Only 46 species have annotated sexes, compared to a range of 125--217 species when not specifying sex.
On this dataset of 46 species, neither the phylogram nor the chronogram is favoured, with undecided support ($0.497 \leq \PhyloSupport \leq 0.848$).
Unfortunately, this result suggests that for a single trait, unless the signal is strong enough, many species are required to obtain a statistically significant conclusion.
}

\DIFadd{More }\DIFaddend generally, our method is also relevant when changes in mutation rate (per generation) and effective population size ($\Ne$) are expected.
First, normalizing by nucleotide divergence also accounts theoretically for variations in effective population size ($\Ne$) since the substitution rate of neutral mutations is equal to the mutation rate, such that $\Ne$ has no effect on the rate of neutral evolution~\citep{kimura_evolutionary_1968, ohta_population_1972}.
Moreover, if population structure were to impact the probability of fixation for neutral mutations (which would also impact a neutral trait), the phylogram would automatically absorb these changes.
This argument was partially tested in our simulation by modelling changes in both $\Ne$ and generation time (although not modelling population structure \textit{per se}), and the phylogram was favoured over the chronogram for a trait evolving neutrally.
Second, under certain conditions, the use of phylograms can also absorb changes in mutation rate (per loci per generation), which in mammals can vary up to 10 fold~\citep{bergeron_evolution_2023}.
As in the case of our simulations, we assumed that the genetic architecture of the trait is constant and the mutation rate is homogeneous across the genome.
However, if the mutation rate would only change for a subset of the genome that is encoding the trait, the phylogram would not absorb these changes in mutation rate.
Finally, because chronograms are usually derived from phylograms by calibrating the tree with molecular clocks, using phylograms directly has the additional advantage of removing model assumptions required to calibrate node ages~\citep{litsios_effects_2012}.

% When does it fall appart, our second result
These examples show that while phylograms are useful to model neutral evolution, they are not a silver bullet to disentangle the effect of drift and selection.
In our simulations, we generally found that using a phylogram tends to favour the OU process over BM.
This suggests that trees in units of nucleotide divergence may not always provide a reliable framework for modelling trait evolution, especially when selection is involved.
We recommend continuing to use chronograms for modelling selection acting on a trait, as they generally provide more accurate and reliable results.
Moreover, for a trait under selection for which the pace of optimum changes is also tracking the changes in generation time, the support for a phylogram over a chronogram would be misleading~\citep{litsios_effects_2012}.
Therefore, we argue that while phylograms can offer valuable insights~\citep{wilson_chronogram_2022}, they should be used with caution and in conjunction with chronograms to provide a more comprehensive understanding of trait evolution.
Contrarily to our cautious note, \citet{wilson_chronogram_2022} suggested that phylograms are more robust for studying discrete character evolution than chronogram.
This apparent contradiction is due to the different units of measurement used in the studies, as they measure branch lengths in units of morphological distance rather than substitutions per site~\citep{wilson_chronogram_2022}.
We thus argue that the use of a phylogram, measured in units of substitutions, could also be useful for studying the evolution of discrete characters, but the same caution should be applied.

% Altogether, what we have done and haven't tested and done yet
Caution is indeed warranted, as there are several limitations to our approach.
First, we simulated only a subset of different regimes of selection\DIFdelbegin \DIFdel{, but }\DIFdelend \DIFaddbegin \DIFadd{.
Notably, we did not simulate a static OU process (with a constant optimum), which would presumably favour an OU model; such simulations are outside the scope of this study since our focus is on distinguishing neutral evolution from selection with a moving optimum.
As emphasized by~}\citet{grabowski_cautionary_2023}\DIFadd{, the main value of OU models lies in multi-optimum variants that can detect shifts between selective regimes.
But }\DIFaddend more complex scenarios could be included, for example with primary and secondary moving optima~\citep{hansen_three_2024}.
\DIFaddbegin 

\DIFaddend Additionally, the simulated genetic architecture of a single trait is assumed to be constant across the phylogeny.
\DIFdelbegin \DIFdel{Constant architecture results in the saturation of the phenotype on long time scales, with a phenomenon similar to nucleotide saturation~}%DIFDELCMD < \citep{latrille_detecting_2024}%%%
\DIFdel{.
This saturation effect might be one of the reasons for a bias in favour of the OU process over BM that we observed in our simulations.
}\DIFdelend A more realistic genetic architecture should include pleiotropy, epistasis and linkage disequilibrium, while here we only considered non-linked loci contributing additively to a single trait.
Second, while the tested alternative models of evolution (multi-rates BM, OU, multi-optimum OU) are standard, they do not encompass the full breadth of available methods to model trait evolution~\citep{pennell_geiger_2014, hohna_revbayes_2016}.
Finally, the mammalian dataset (brain and body) may not be representative of all traits or species and even though our analysis showcases the utility of phylograms, more empirical studies are needed to replicate the findings.

% Gene expression level and empirical studies.
\DIFdelbegin \DIFdel{One striking example of a trait that could (and should) be empirically studied with a phylogram is gene expression level, both for }\DIFdelend \DIFaddbegin \DIFadd{Gene expression levels (both }\DIFaddend mRNA and protein levels\DIFdelbegin \DIFdel{.
}\DIFdelend \DIFaddbegin \DIFadd{) represent one class of traits where distinguishing neutral evolution from selection is particularly important.
While the polygenic nature of expression traits can vary, testing whether their evolution is better explained by phylograms or chronograms can help determine whether changes are driven by accumulation of neutral substitutions or by selection responding to environmental changes.
}\DIFaddend Indeed, the regime of selection acting on gene expression level is the focus of intense debate~\citep{signor_evolution_2018, price_detecting_2022, bertram_cagee_2023, dimayacyac_evaluating_2023}.
Typically, datasets on which BM is favoured over other models is interpreted as mRNA expression level evolving neutrally~\citep{khaitovich_evolution_2006, catalan_drift_2019, dimayacyac_evaluating_2023}.
Phenotypic distance is also plotted as a function of time while claiming that a trait is evolving under drift~\citep{jiang_decoupling_2023}.
Instead, we argue that phenotypic distance should be \DIFdelbegin \DIFdel{shown as a function of branch length }\DIFdelend \DIFaddbegin \DIFadd{expressed }\DIFaddend in units of substitutions \DIFdelbegin \DIFdel{(square root ), not time}\DIFdelend \DIFaddbegin \DIFadd{instead of time, using the square root of branch lengths}\DIFaddend .
More generally, from an empirical perspective, the use of nucleotide changes at neutral loci as a normalizing factor brings new perspectives into trait evolution~\citep{latrille_detecting_2024}.

% Conclusion
Altogether, our study supports the use of a chronogram when testing for selection acting on a trait.
\DIFdelbegin \DIFdel{The pipeline of hypothesis testing and model comparison is currently well established and should be followed.
However, when the analysis concludes that the BM is favoured, then the support of the phylogram over the chronogram should be tested additionally}\DIFdelend \DIFaddbegin \DIFadd{We propose the phylogram-chronogram comparison as a new tool in the comparative methods toolbox, particularly useful when seeking deviation from neutral evolution.
In such cases, comparing the support for a phylogram versus a chronogram provides an additional diagnostic to assess the regime of evolution acting on a trait}\DIFaddend .
We provide a Bayesian model implemented in \textit{RevBayes} to \DIFdelbegin \DIFdel{test this hypothesis, which indicates whether a BM supports better }\DIFdelend \DIFaddbegin \DIFadd{perform this comparison, which estimates whether BM supports }\DIFaddend a phylogram or a chronogram ($0 \leq \PhyloSupport \leq 1$), given the data at the tip of the tree and both trees with the same topology.
\DIFaddbegin \DIFadd{For empiricists more familiar with frequentist approaches, we also provide a maximum likelihood implementation that yields comparable results through AIC weights (see supplementary material section~\ref{sec:maximum-likelihood-estimation}).
}\DIFaddend Only if the phylogram has more support than the chronogram ($\PhyloSupport$ close to $1$) \DIFdelbegin \DIFdel{can it be assumed that the trait is evolving under drift, but caution is still warranted}\DIFdelend \DIFaddbegin \DIFadd{should one consider the neutral model of evolution as a plausible explanation for the evolution of the trait, but claims of neutral evolution should still be made cautiously}\DIFaddend .

%TC:ignore

\section*{Author contributions}
Original idea: T.L.\ and anonymous reviewers;
Model conception: T.L., T.G.\ and N.S.;
Code: T.L.;
Data analyses: T.L.\ and T.G.;
Interpretation: T.L., T.G.\ and N.S.;
First draft: T.L.;
Editing and revisions: T.L., T.G.\ and N.S.
Project management and funding: N.S\@.

\section*{Acknowledgements}
\label{sec:acknowledgment}
We gratefully acknowledge the help of Lucy M. Fitzgerald, Diego A. Hartasánchez, Anna Marcionetti and Julien Joseph for their advice and reviews concerning this manuscript.

\section*{Funding}
This work was funded by Faculté de Biologie et de Médecine, Université de Lausanne (https://www.unil.ch; to TL, TG and NS) and Swiss National Science Fund (https://www.snf.ch; grant 315230-219757 to TL and NS). The funders did not play any role in the study design, data collection and analysis, decision to publish, or preparation of the manuscript.

\section*{Competing interests}
The authors declare no conflicts of interest.

\subsection*{Data and materials availability:}
The materials that support the findings of this study are openly available in GitHub at \href{https://github.com/ThibaultLatrille/ChronoPhylogram}{github.com/ThibaultLatrille/ChronoPhylogram}.
Snakemake pipeline, simulator, analysis scripts and \textit{RevBayes} scripts and documentation are available in the repository to replicate the study.

\printbibliography

\newpage

\part*{Supplementary materials}
\renewcommand{\thetable}{S\arabic{table}}
\renewcommand{\thefigure}{S\arabic{figure}}
\setcounter{figure}{0}
\setcounter{table}{0}
\setcounter{section}{0}

\renewcommand{\baselinestretch}{1.0}\normalsize
\tableofcontents
\DIFdelbegin %DIFDELCMD < \renewcommand{\baselinestretch}{1.5}%%%
\DIFdelend \DIFaddbegin \renewcommand{\baselinestretch}{1.2}\DIFaddend \normalsize

\DIFaddbegin \newpage

\section{\DIFadd{Theoretical rationale}}\label{sec:theoretical-rationale}

\DIFadd{In this section, we derive the theoretical expectation for the relationship between phenotypic covariance and nucleotide divergence under neutral evolution, using the infinitesimal model of quantitative genetics~}\citep{lande_quantitative_1979, lande_sexual_1980} \DIFadd{and notations as from }\citet{latrille_detecting_2024}\DIFadd{.
}

\subsection{\DIFadd{Genetic architecture and mutational variance}}

\DIFadd{For a given quantitative trait, the genetic architecture is mainly defined by the number of loci encoding the trait ($\NbrLoci$) and the random additive effect of a mutation on the trait ($a$).
For a diploid individual, the mutational variance ($\VarMutation$) is the rate at which trait variance changes between generations due to mutations.
As shown in }\citet{lande_quantitative_1979, lande_sexual_1980}\DIFadd{, $\VarMutation$ is a function of the mutation rate per locus per generation ($\MutationRatePheno$) and the genetic architecture of the trait:
}\begin{align}
    \DIFadd{\VarMutation = 2 \MutationRatePheno \Multiply \GenArchi \label{eq:var-mutation}.
}\end{align}

\DIFadd{Under the infinitesimal model, mutations supply new genetic variants while random genetic drift depletes standing variation~}\citep{turelli_commentary_2017, barton_infinitesimal_2017, sella_thinking_2019}\DIFadd{.
For a neutral trait at mutation-drift equilibrium~}\citep{lynch_mutation_1998}\DIFadd{, the additive genetic variance ($\VarGenetic$) is a function of the mutational variance ($\VarMutation$) and the effective population size ($\Ne$):
}\begin{align}
    \DIFadd{\VarGenetic }& \DIFadd{= 2 \Ne \Multiply \VarMutation, }\\
    & \DIFadd{= 4 \Ne \Multiply \MutationRatePheno \Multiply \GenArchi \text{ from eq.~\ref{eq:var-mutation}}. \label{eq:var-genetic}
}\end{align}

\subsection{\DIFadd{Covariance between species}}

\DIFadd{For a given species $\Spi$, we denote by $\MeanTrait_{\Spi}$ the mean value of the trait across the individuals.
If the trait is neutral and encoded by many loci as assumed by the infinitesimal model, $\MeanTrait_{\Spi}$ evolves as a Brownian process~}\citep{felsenstein_phylogenies_1985, hansen_translating_1996}\DIFadd{.
}

\DIFadd{Given a phylogenetic tree, for a pair of species $\Spi$ and $\Spj$ from this tree, we denote as $\NbrGen$ the number of generations between the root of the tree and the most recent common ancestor of species $\Spi$ and $\Spj$.
Crucially, the number of generations depends on both the elapsed time ($\Time_{\Spi, \Spj}$) and the generation time ($\GenerationTime$, the average time between two consecutive generations):
}\begin{align}
    \DIFadd{\NbrGen = \frac{\Time_{\Spi, \Spj}}{\GenerationTime}. \label{eq:gen-time}
}\end{align}

\DIFadd{The covariance between $\MeanTrait_{\Spi}$ and $\MeanTrait_{\Spj}$ depends on the additive genetic variance and the number of shared generations, as given by }\citet{hansen_translating_1996}\DIFadd{:
}\begin{align}
    \DIFadd{\VarPhy }& \DIFadd{= \frac{\VarGenetic}{\Ne} \Multiply \NbrGen, }\\
    & \DIFadd{= 4 \NbrGen \Multiply \MutationRatePheno \Multiply \GenArchi \text{ from eq.~\ref{eq:var-genetic}}. \label{eq:covar}
}\end{align}

\DIFadd{Substituting eq.~\ref{eq:gen-time} into eq.~\ref{eq:covar}, we can express the covariance in terms of time:
}\begin{align}
    \DIFadd{\VarPhy }& \DIFadd{= 4 \Multiply \frac{\Time_{\Spi, \Spj}}{\GenerationTime} \Multiply \MutationRatePheno \Multiply \GenArchi. \label{eq:covar-time}
}\end{align}
\DIFadd{Equation~\ref{eq:covar-time} shows that trait covariance depends not only on elapsed time but also inversely on generation time: lineages with shorter generation times accumulate more generations per unit time and thus greater phenotypic divergence.
This dependence on generation time is problematic when using chronograms, as generation time can vary substantially across lineages~}\citep{demagalhaes_database_2009}\DIFadd{.
}

\subsection{\DIFadd{Neutral substitution rate}}

\DIFadd{For any genomic region under neutral evolution, some mutations will eventually reach fixation due to random genetic drift, resulting in a nucleotide substitution at the species level.
The probability of fixation ($\pfix$) of a neutral mutation in a diploid population is $1/2\Ne$~}\citep{kimura_probability_1962}\DIFadd{.
We can derive the substitution rate ($\SubRate$, per nucleotide site per generation) as the number of newly arisen mutations ($2\Ne \Multiply \MutationRateNuc$, where $\MutationRateNuc$ is the mutation rate per nucleotide site per generation) multiplied by the probability of fixation~}\citep{kimura_evolutionary_1968}\DIFadd{:
}\begin{align}
    \DIFadd{\SubRate }& \DIFadd{= 2 \Ne \Multiply \MutationRateNuc \Multiply \pfix, }\\
    & \DIFadd{= 2 \Ne \Multiply \MutationRateNuc \Multiply \dfrac{1}{2\Ne}, }\\
    & \DIFadd{= \MutationRateNuc. \label{eq:substitution-rate}
}\end{align}
\DIFadd{This fundamental result shows that for neutral mutations, the substitution rate per generation equals the mutation rate per generation, independent of population size~}\citep{kimura_evolutionary_1968, mccandlish_modeling_2014}\DIFadd{.
}

\subsection{\DIFadd{Nucleotide divergence absorbs generation time}}

\DIFadd{We denote $\NucDiv$ as the nucleotide divergence between the root of the tree and the most recent common ancestor of species $\Spi$ and $\Spj$.
In other words, $\NucDiv$ is the expected number of substitutions per nucleotide site during $\NbrGen$ generations.
Assuming that no multiple substitutions occurred at the same site (i.e., low saturation), $\NucDiv$ is:
}\begin{align}
    \DIFadd{\NucDiv }& \DIFadd{= \NbrGen \Multiply \SubRate, }\\
    & \DIFadd{= \NbrGen \Multiply \MutationRateNuc \text{ from eq.~\ref{eq:substitution-rate}}. \label{eq:distance}
}\end{align}

\DIFadd{From eq.~\ref{eq:distance}, we can express the number of generations as:
}\begin{align}
    \DIFadd{\NbrGen = \frac{\NucDiv}{\MutationRateNuc}. \label{eq:gen-from-div}
}\end{align}

\subsection{\DIFadd{Trait covariance scales with nucleotide divergence}}

\DIFadd{Substituting eq.~\ref{eq:gen-from-div} into eq.~\ref{eq:covar}, we obtain:
}\begin{align}
    \DIFadd{\VarPhy }& \DIFadd{= 4 \Multiply \frac{\NucDiv}{\MutationRateNuc} \Multiply \MutationRatePheno \Multiply \GenArchi, }\\
    & \DIFadd{= \NucDiv \Multiply \frac{\MutationRatePheno}{\MutationRateNuc} \Multiply 4\GenArchi . \label{eq:covar-div}
}\end{align}

\DIFadd{Equation~\ref{eq:covar-div} shows that the covariance in mean trait value between a pair of species ($\VarPhy$) increases linearly with nucleotide divergence ($\NucDiv$), with a proportionality constant that depends on the genetic architecture of the trait and the ratio of phenotypic to nucleotide mutation rates.
This derivation shows first that by expressing trait covariance in terms of nucleotide divergence rather than time (compare eq.~\ref{eq:covar-time} with eq.~\ref{eq:covar-div}), the generation time ($\GenerationTime$) is removed from the equation.
Lineages with different generation times will have proportionally different nucleotide divergences, automatically accounting for the variation in the number of generations.
Second, the effective population size ($\Ne$) appears in both the additive genetic variance (eq.~\ref{eq:var-genetic}) and implicitly in the substitution rate derivation, but cancels when considering neutral trait and nucleotide divergence.
}

\DIFadd{Concerning changes in mutation rates, the proportionality constant in eq.~\ref{eq:covar-div} depends on the ratio of phenotypic to nucleotide mutation rates ($\MutationRatePheno / \MutationRateNuc$).
If changes in mutation rate affect the entire genome proportionally (including both neutral loci and trait-encoding loci), this ratio remains constant and mutation rate (per generation) variation across lineages is also absorbed.
Importantly, in vertebrates the mutation rate per generation is varying less than the mutation rate per unit of time~}\citep{bergeron_evolution_2023}\DIFadd{.
Finally, the derivation assumes that the genetic architecture ($\GenArchi$) remains constant across the phylogeny.
Changes in the number of loci or the distribution of mutational effects would violate this assumption.
}

\newpage
\DIFaddend \section{Simulation of continuous trait}\label{sec:simulator}

\subsection{Genotype-phenotype map}\label{subsec:genotype-phenotype-map}

\begin{itemize}
    \item $\NbrLoci$ is the number of loci encoding the trait.
    \item $a_l \sim \mathcal{N}(0,a^2)$ is the effect of a mutation on the trait at locus $l \in \{1, \hdots, \NbrLoci\}$.
    \item $\Ne$ is the effective number of individuals.
    \item $g_{\Indiv,l} \in \{0, 1, 2\}$ is the genotypic value at locus $l$ for individual $\Indiv \in \{1, \hdots, \Ne\}$.
    \item $G_{\Indiv} = \sum_{l=1}^{\NbrLoci} a_l \times g_{\Indiv,l}$ is the genotypic value for individual $\Indiv$.
    \item $\xi_{\Indiv} \sim \mathcal{N}(0, \VarEnv)$ is the effect of environment on the trait for individual $\Indiv$.
    \item $\Trait_{\Indiv} = G_{\Indiv}+ \xi_i$ is the phenotype for individual $\Indiv$.
\end{itemize}

\begin{center}
    \captionof{figure}{summary of trait's genetic architecture.}
    \includegraphics[width=0.7\textwidth, page=1] {figures/artwork_genet_architecture}
    \label{fig:simulator-summary}
\end{center}

\subsection{Simulation along a tree}\label{subsec:simulation-along-a-tree}
% vG: 13.8, vP: 69.0, vE: 55.2, nbr_loci: 5,000, a: 1.0, mut_rate: 1.38e-05, pop_size: 50, h2: 0.2
% pS = 0.0027577478571428568
% Tree length (dS) = 2.340579
% Tree length (My) = 1259.4490200000002
% Root age (My) = 98.9489413888889
% Root age (dS) = 0.18388820079995308
% nbr_sites_var = 27.577478571428568
% u = 1.3788739285714283e-05
% nbr_generations = 13336.114128321291
Each individual phenotypic value was the sum of genotypic value and an environmental effect.
The environmental effect was normally distributed with variance $\VarEnv$.
We assumed that the genotypic value was encoded by $\NbrLoci=5,000$ loci, with each locus contributing an additive effect that was normally distributed with standard deviation $a=1$.
\DIFaddbegin \DIFadd{Narrow-sense heritability is defined as $\Heritability = \VarGenetic / (\VarGenetic + \VarEnv)$, where $\VarGenetic$ is the additive genetic variance.
}\DIFaddend We assumed a trait with a narrow-sense heritability of $\Heritability=0.2$ \DIFaddbegin \DIFadd{at the root of the tree }\DIFaddend and computed the theoretical $\VarEnv$ accordingly\DIFdelbegin \DIFdel{.
}\DIFdelend \DIFaddbegin \DIFadd{, using the formula $\VarEnv = \VarGenetic \Multiply (1 - \Heritability) / \Heritability$, with $\VarGenetic$ calculated as in eq.~\ref{eq:var-genetic} from the genetic architecture of the trait, the mutation rate per locus per generation ($\MutationRatePheno$) and the effective population size ($\Ne$) at the root of the tree.
We held the $\VarEnv$ constant across the tree, which means that in practice the heritability of the trait can vary across lineages as $\VarGenetic$ changes (due to changes in $\Ne$ and $\MutationRatePheno$).
}

\DIFaddend Assuming a diploid panmictic population of size $\Ne=50$ at the root of the tree, and with non-overlapping generations, we simulated explicitly each generation along an ultrametric phylogenetic tree.
For each offspring, the number of mutations was drawn from a Poisson distribution with mean $2 \Multiply \MutationRatePheno \Multiply \NbrLoci $, with the mutation rate per locus per generation $\MutationRatePheno$.
\DIFaddbegin 

\DIFaddend From the empirical mammalian dataset (see Methods), we computed an average nucleotide divergence from the root to leaves of $0.18$.
We scaled parameters in our simulations to fit plausible values for mammals.
We thus used a nucleotide mutation rate of $\MutationRateNuc=0.00276 / 4 \Ne = 1.38 \times 10^{-5}$ per site per generation and a total of $0.18 / 1.38 \times 10^{-5} = 13,500$ generations from root to leaves, and the number of generations along each branch was proportional to the branch length.
We set $\MutationRatePheno=\MutationRateNuc$ without loss in generality since the genetic architecture ($\NbrLoci$ and $a$) is assumed constant in the simulator.

\DIFaddbegin \newpage
\DIFaddend The changes in $\MutationRatePheno$ and $\Ne$ along the lineages were both modelled by a Brownian motion (BM) on the log scale (log-$\MutationRatePheno$ and log-$\Ne$), leading to geometric Brownian motion on the linear scale ($\MutationRatePheno$ and $\Ne$).
These processes are parametrized as $\brownian \left(0, \sigma_{\MutationRatePheno}=0.0043\right)$ and $\brownian \left(0, \sigma_{\Ne}=0.0043\right)$, which, if counted across $13,500$ generations, leads to a standard deviation of $0.0043 \Multiply \sqrt {13,500} = 0.5$.
In other words, the deviation in log-$\Ne$ and log-$\MutationRatePheno$  between the extant species and the root is $0.5$.
\DIFaddbegin 

\DIFaddend An Ornstein-Uhlenbeck process was overlaid to the instant value of log-$\Ne$ provided by the geometric BM to account for short-term changes between generations ($\text{OU} \left(0, \sigma_{\Ne}=0.1, \theta_{\Ne}=0.9\right)$).
The geometric Brownian motion accounted for long-term fluctuations (low rate of changes $\sigma_{\Ne}$ but unbounded), while the Ornstein-Uhlenbeck introduced short-term fluctuations (high rate of changes $\sigma_{\Ne}$ but bounded and mean-reverting).
The simulation started from an initial sequence at equilibrium at the root of the tree and, at each node, the process was split until it finally reached the leaves of the tree.
From a speciation process perspective, this was equivalent to an allopatric speciation over one generation.

At each generation, parents were randomly sampled with a weight proportional to their fitness ($W$).
Selection was modelled as a one-dimensional Fisher's geometric landscape, with the fitness of an individual being a \DIFdelbegin \DIFdel{monotonously }\DIFdelend \DIFaddbegin \DIFadd{monotonically }\DIFaddend decreasing function of the distance between the individual and the optimal phenotype~\citep{tenaillon_utility_2014,blanquart_epistasis_2016}.
More specifically, the fitness of an individual was given by $W = \e^{(\Trait - \lambda)^2/ \alpha}$, where $\Trait$ was the trait value of the individual, $\lambda=0.0$ was the optimal trait value, and $\alpha=0.02$ was the strength of selection.
Mutations were considered as a displacement of the phenotype in the multidimensional space.
Beneficial mutations moved the phenotype closer to the optimum, while deleterious mutations moved it further away.
\DIFaddbegin 

\DIFaddend Moving optimum selection was implemented by allowing the optimum phenotype to move along the phylogenetic tree as a geometric BM~\citep{hansen_stabilizing_1997} ($\lambda \sim \brownian \left(0, \sigma_{\lambda}=1.0\right)$).
Multiple optima were implemented by allowing the optimum phenotype to shift along a branch with a probability of $10^{-5}$, the shift in optimum being then drawn from a exponential distribution with rate parameter $0.1$, meaning an average of $10$ in the optimum shift per switch.
Neutral evolution was implemented by flattening the fitness landscape ($W=1$), which meant that each individual had the same probability of being sampled at each generation.

\newpage
\section{Ancestral trait reconstruction}\label{sec:ancestral-trait-reconstruction}
\begin{center}
    \begin{minipage}{0.4\linewidth}
        \flushleft {\tiny  \textbf{A}}
        \includegraphics[width=\linewidth, page=1]{figures/simulation_varGT_dist_NeutralPhylo}
    \end{minipage}
    \begin{minipage}{0.4\linewidth}
        \flushleft {\tiny  \textbf{B}}
        \includegraphics[width=\linewidth, page=1]{figures/simulation_varGT_dist_NeutralChrono}
    \end{minipage}\\
    \begin{minipage}{0.4\linewidth}
        \flushleft {\tiny  \textbf{C}}
        \includegraphics[width=\linewidth, page=1]{figures/simulation_varGT_dist_SelectionPhylo}
    \end{minipage}
    \begin{minipage}{0.4\linewidth}
        \flushleft {\tiny  \textbf{D}}
        \includegraphics[width=\linewidth, page=1]{figures/simulation_varGT_dist_SelectionChrono}
    \end{minipage}
    \captionof{figure}{
        \textbf{Phenotypic distance as a function of branch length.}
        Phenotypic distance ($|\Delta z|$) is computed as the absolute difference of the posterior mean trait value between the ancestral node and the descendent node.
        Across all simulated dataset (100 replicates), we computed the mean phenotypic distance (blue dots) and standard deviation (blue horizontal lines) for each branch.
        Simulations under a neutral regime (yellow, top row, panels A and B) and under a moving optimum regime (bottom row, panels C and D).
        Reconstruction of ancestral trait on a phylogram (blue, left column, panels A and C) and a chronogram (right column, panels B and D).
        \label{fig:results-distance}
    }
\end{center}

\newpage
\section{Empirical dataset}\label{sec:empirical-dataset}

\begin{center}
\begin{adjustbox}{width = 0.75\textwidth}
\begin{tabular}{||l|l|l|l|r|r||}
\hline \textbf{Dataset} & \textbf{Trait} & \textbf{Sex} & \textbf{Likelihood} & \textbf{Taxa} & \textbf{Phylogram support ($\bm{\pi}$)} \\
\hline \hline
\textit{Mammal} & Body mass & Mixed & Full & 125 & 0.0 \\ \hline
\textit{Mammal} & Body mass & Mixed & REML & 125 & 0.002 \\ \hline
\textit{Mammal} & Body mass & \Female & Full & 46 & 0.848 \\ \hline
\textit{Mammal} & Body mass & \Female & REML & 46 & 0.846 \\ \hline
\textit{Mammal} & Body mass & \Male & Full & 46 & 0.697 \\ \hline
\textit{Mammal} & Body mass & \Male & REML & 46 & 0.684 \\ \hline
\textit{Mammal} & Brain mass & Mixed & Full & 125 & 0.0035 \\ \hline
\textit{Mammal} & Brain mass & Mixed & REML & 125 & 0.0035 \\ \hline
\textit{Mammal} & Brain mass & \Female & Full & 46 & 0.56 \\ \hline
\textit{Mammal} & Brain mass & \Female & REML & 46 & 0.539 \\ \hline
\textit{Mammal} & Brain mass & \Male & Full & 46 & 0.515 \\ \hline
\textit{Mammal} & Brain mass & \Male & REML & 46 & 0.497 \\ \hline
\textit{COMBINE} & Body mass & Mixed & Full & 217 & 0.0 \\ \hline
\textit{COMBINE} & Body mass & Mixed & REML & 217 & 0.0 \\ \hline
\textit{COMBINE} & Brain mass & Mixed & Full & 174 & 0.0015 \\ \hline
\textit{COMBINE} & Brain mass & Mixed & REML & 174 & 0.00649 \\ \hline
\end{tabular}
\end{adjustbox}
\DIFdelbegin %DIFDELCMD < \captionof{table}{
%DIFDELCMD < \textbf{Summary of empirical datasets.}
%DIFDELCMD < Support for the phylogram over the chronogram ($\pi$) for different datasets and traits.
%DIFDELCMD < Brain and body mass is either extracted from \citet[\textit{\textit{Mammal}}]{tsuboi_breakdown_2018} or from the \textit{COMBINE} database~\citep{soria_combine_2021}.
%DIFDELCMD < For the \citet[\textit{\textit{Mammal}}]{tsuboi_breakdown_2018} dataset, the individuals can also be split into males (\Male) and females (\Female).
%DIFDELCMD < Likelihood is either computed with a full likelihood or with a restricted maximum likelihood (REML).
%DIFDELCMD < The number of taxa is the number of species with data available for the trait.
%DIFDELCMD < \label{table:results-empirical}
%DIFDELCMD < }
%DIFDELCMD < %%%
\DIFdelend \DIFaddbegin \captionof{table}{
\textbf{Summary of empirical datasets.}
Support for the phylogram over the chronogram ($\pi$) for different datasets and traits.
Brain and body mass is either extracted from \citet[\textit{\textit{Mammal}}]{tsuboi_breakdown_2018} or from the \textit{COMBINE} database~\citep{soria_combine_2021}.
For the \citet[\textit{\textit{Mammal}}]{tsuboi_breakdown_2018} dataset, the individuals can also be split into males (\Male) and females (\Female).
Likelihood is either computed with a full likelihood or with a restricted maximum likelihood (REML).
The number of species is the number of species with data available for the trait.
\label{table:results-empirical}
}
\end{center}

\newpage
\section{\DIFadd{Maximum Likelihood estimation}}\label{sec:maximum-likelihood-estimation}

\subsubsection*{\DIFadd{Maximum likelihood estimate}}
\DIFadd{At the phylogenetic scale, for $\NbrTaxa$ species in the tree, $\DistanceMatrix$ ($\NbrTaxa \times \NbrTaxa$) is the symmetric distance matrix computed from the branch lengths and the topology of the phylogenetic tree.
The diagonal elements $\Distance_{\Spi,\Spi}$ represents the total distance (sum of branch length) from the root of the tree to each species ($\Spi$).
The off-diagonal elements of the distance matrix between a pair of species ($\Distance_{\Spi,\Spj} = \NucDiv$) are the distances between the root and the most recent common ancestor (MRCA) of species $\Spi$ and $\Spj$, computed as the sum of branch lengths from the root to the MRCA.
}

\DIFadd{The mean trait value at the root of the tree ($\EstRootTrait$) can be estimated from the $\NbrTaxa \times 1$ vector of mean trait values $\VecTrait$ at the tips of the tree using maximum likelihood~}\citep{omeara_testing_2006}\DIFadd{:
%DIF >  invC = np.linalg.inv(C)
%DIF >  ones = np.ones(n)
%DIF >  v = np.dot(ones.T, invC)
%DIF >  anc_z = float(np.dot(v, x)) / float(np.dot(v, ones))
}\begin{gather}
    \DIFadd{\EstRootTrait = \left( \VecOne\tr \MultiplyMatrix \DistanceMatrix\inv \MultiplyMatrix \VecOne \right)\inv \Multiply \left( \VecOne\tr \MultiplyMatrix \DistanceMatrix\inv \MultiplyMatrix \VecTrait \right), \label{eq:estimated-root-trait}
}\end{gather}
\DIFadd{where $\VecOne$ is an $\NbrTaxa \times 1$ column vector of ones.
}

\DIFadd{The maximum likelihood estimate of the Brownian rate ($\RateBrownian$) can then be computed as~}\citep{omeara_testing_2006}\DIFadd{:
%DIF >  d = (x - anc_z * ones)
%DIF >  var = float(np.dot(d.T, np.dot(invC, d)) / (n - 1))
}\begin{gather}
    \DIFadd{\EstRateBrownian = \frac{1}{\NbrTaxa - 1} \Multiply  \left( \VecTrait - \EstRootTrait \Multiply \VecOne \right)\tr \MultiplyMatrix \DistanceMatrix\inv \MultiplyMatrix \left( \VecTrait - \EstRootTrait \Multiply \VecOne \right). \label{eq:estimated-brownian-rate}
}\end{gather}
\DIFadd{Then, from the estimated paramaters $\EstRootTrait$ and $\EstRateBrownian$ the likelihood ($\mathcal{L}$) of the data ($\VecTrait$) given the tree ($\DistanceMatrix$) and the parameters can be computed as~}\citep{omeara_testing_2006, harmon_phylogenetic_2018}\DIFadd{:
}\begin{gather}
    \DIFadd{\mathcal{L} \left(\VecTrait \ | \ \DistanceMatrix, \EstRootTrait, \EstRateBrownian \right) = \frac{\e^{-1/2 (\VecTrait-\EstRootTrait \Multiply \VecOne)\tr \MultiplyMatrix (\EstRateBrownian \Multiply \DistanceMatrix)\inv \MultiplyMatrix (\VecTrait-\EstRootTrait \Multiply \VecOne)}}{\sqrt{(2 \pi)^\NbrTaxa \Multiply \det (\RateBrownian \Multiply \DistanceMatrix) }} \label{eq:likelihood}
}\end{gather}
\DIFadd{Finally, the log-likelihood of the data ($\ln \mathcal{L}$) is:
%DIF >   ll = -0.5 * (n * np.log(2 * np.pi * var) + np.log(np.linalg.det(C)) + (1 / var) * np.dot(d.T, np.dot(invC, d)))
}\begin{align}
    \DIFadd{\ln \mathcal{L} \left(\VecTrait \ | \ \DistanceMatrix, \EstRootTrait, \EstRateBrownian \right) }&\DIFadd{= \ln \left( \frac{\e^{-1/2 (\VecTrait-\EstRootTrait \Multiply \VecOne)\tr \MultiplyMatrix (\EstRateBrownian \Multiply \DistanceMatrix)\inv \MultiplyMatrix (\VecTrait-\EstRootTrait \Multiply \VecOne)}}{\sqrt{(2 \pi)^\NbrTaxa \Multiply \det (\RateBrownian \Multiply \DistanceMatrix) }} \right)\text{ from eq.~\ref{eq:likelihood}}, }\\
    &\DIFadd{= -\frac{1}{2} \left[ (\VecTrait-\EstRootTrait \Multiply \VecOne)\tr \MultiplyMatrix (\EstRateBrownian \Multiply \DistanceMatrix)\inv \MultiplyMatrix (\VecTrait-\EstRootTrait \Multiply \VecOne) \right] - \ln \left( \sqrt{ (2 \pi)^\NbrTaxa \Multiply \det (\RateBrownian \Multiply \DistanceMatrix) } \right) , }\\
    &\DIFadd{= -\frac{1}{2} \left[\frac{1}{\EstRateBrownian } \Multiply (\VecTrait-\EstRootTrait \Multiply \VecOne)\tr \MultiplyMatrix  \DistanceMatrix\inv \MultiplyMatrix (\VecTrait-\EstRootTrait \Multiply \VecOne) + \ln \left( (2 \pi)^\NbrTaxa \Multiply \det (\RateBrownian \Multiply \DistanceMatrix)  \right) \right], }\\
    &\DIFadd{= -\frac{1}{2} \left[\NbrTaxa - 1  + \ln \left( (2 \pi)^\NbrTaxa \Multiply (\RateBrownian)^\NbrTaxa \Multiply \det ( \DistanceMatrix)  \right) \right]\text{ from eq.~\ref{eq:estimated-brownian-rate}}, }\\
    &\DIFadd{= -\frac{1}{2} \Multiply \left [\NbrTaxa \Multiply \ln \left[2\pi \Multiply \EstRateBrownian\right] + \ln \left[\det \left(\DistanceMatrix\right) \right] + \NbrTaxa - 1 \right]. \label{eq:log-likelihood}
}\end{align}

\subsubsection*{\DIFadd{AIC weights for model comparison}}
\DIFadd{To compare the fit of BM on the phylogram versus the chronogram, we compute the log-likelihood for each tree type using equations~\ref{eq:estimated-root-trait}--\ref{eq:log-likelihood}.
The Akaike Information Criterion (AIC) for each model is:
}\begin{gather}
    \DIFadd{\text{AIC} = 2k - 2 \ln \mathcal{L},
}\end{gather}
\DIFadd{where $k=2$ is the number of parameters (root trait value and Brownian rate).
The relative likelihood of each model is computed as:
}\begin{gather}
    \DIFadd{w_{\text{phylo}} = \frac{\e^{-\Delta_{\text{phylo}}/2}}{\e^{-\Delta_{\text{phylo}}/2} + \e^{-\Delta_{\text{chrono}}/2}},
}\end{gather}
\DIFadd{where $\Delta_{\text{phylo}} = \text{AIC}_{\text{phylo}} - \min(\text{AIC}_{\text{phylo}}, \text{AIC}_{\text{chrono}})$ and similarly for $\Delta_{\text{chrono}}$.
The AIC weight $w_{\text{phylo}}$ is analogous to the posterior mean $\PhyloSupport$ in the Bayesian framework.
}

\newpage
\begin{center}
    \begin{minipage}{0.49\linewidth}
        \flushleft {\tiny  \textbf{\DIFadd{A}}}
        \includegraphics[width=\linewidth, page=1]{figures/simulation_varGT_ML_vs_Bayes}
    \end{minipage}
    \begin{minipage}{0.49\linewidth}
        \flushleft {\tiny  \textbf{\DIFadd{B}}}
        \includegraphics[width=\linewidth, page=1]{figures/simulation_varGT_ML_vs_Bayes_logit}
    \end{minipage}
    \captionof{figure}{
        \textbf{Comparison of Bayesian and frequentist approaches.}
        Posterior mean of $\PhyloSupport$ (Bayesian approach) versus AIC weight for the phylogram (frequentist approach) across 100 replicate simulations.
        Points are coloured by the regime of evolution: neutral (yellow), moving optimum (blue), and multiple optima (red).
        The strong correlation between the two metrics indicates that both approaches yield consistent conclusions about the relative support for phylogram versus chronogram.
        Panel~A: direct comparison.
        Panel~B: comparison on the logit scale, to better visualize extreme values close to 0 and 1. By definition, $\logit(x) = \ln (x/(1-x))$.
        \label{fig:aicweights}
    }
\DIFaddend \end{center}

%TC:endignore
\end{document}